\documentclass[11pt]{article}
\usepackage{amsmath,amssymb,amsthm,mathtools,hyperref,geometry}
\geometry{margin=1in}
\newtheorem{theorem}{Theorem}
\newtheorem{lemma}[theorem]{Lemma}
\newtheorem{proposition}[theorem]{Proposition}
\newtheorem{corollary}[theorem]{Corollary}
\newcommand{\Set}{\mathbf{Set}}
\newcommand{\Adm}{\operatorname{Adm}}
\newcommand{\Coup}{\operatorname{Coup}}
\newcommand{\id}{\mathrm{id}}
\title{Local Joint Comparisons, Native Grade Lifting and Generated Interference}
\date{15 November 2026}
\begin{document}
\maketitle

\paragraph{Claim and boundary.} Primitive native comparisons extend to finite executions generated by the rule format below when local joint admission is closed under comparison transport, and a fixed joint chooser descends through the retained input comparison. The whole-system joint kernel square is a conclusion. The initial fragment uses monadic effect-fibre contexts and one-hole writer content, with actual shared-state microsteps. Section~\ref{sec:grade-lifting} constructs native grade lifts from independent coefficient limits, proves their rule coherence and recovers an original augmented-vector-space doctrine with nonfaithful extensions. It preserves directed cells, every supplied coupling choice, complete raw witnesses and nonserializable histories. It is not a theorem for arbitrary contextad/contentad pairs, arbitrary physical coupling, instantaneous feedback, or unrestricted infinite-effect completion. Section~\ref{sec:graded-feedback} derives shared-state closure and augmented native history-prefix systems, and separates their interpretation from infinite-effect and infinite-grade completion. Section~\ref{sec:occurrence-realization} supplies the second asynchronous realization; Section~\ref{sec:core-completion-assessment} assesses the resulting scoped v0.1.3 completion evidence and remaining programme obligations. The local transport lemma combines existing tree admission and native cone results; no historical novelty is asserted.

The question addressed is a design-and-proof question: which local law makes native comparisons closed under generated interference? The governing observations include the full retained body, event chronology, shared world, enabled tasks and local information. The construction identifies only canonical tuple bracketings. It does not identify couplings with equal marginals, histories with equal final states, zero-weight witnesses with absence, or different native grade maps with identical carrier maps. Its source-level interpretation remains conditional on an adequate presentation of primitive effects and the chosen interference boundary.

\section{Native effects and local joint data}

Fix a monad $T$ on $\Set$. For a step with input $a\in A$, supply a raw witness set $W_a$, an effect $p_a\in TW_a$, and a full body
\[
 f_a:W_a\longrightarrow B\times Q.
\]
A comparison to a primed step has pure maps $h:A\to A'$, $k:B\to B'$, a content map $v:Q\to Q'$, and actual maps $r_a:W_a\to W'_{ha}$ such that
\[
 Tr_a(p_a)=p'_{ha},\qquad f'_{ha}r_a=(k\times v)f_a. \tag{N}\label{eq:native}
\]
No map is presumed invertible. Raw carriers include declared zero-weight cases. $Q$ can retain an event word, receipt, cost, pending obligation or other task-relevant output; omitting such a coordinate changes the claim.

At an input $x$, take a finite nonempty overlap tree with vertices $i$, witness sets $W_{i,x}$, effects $p_{i,x}$, and boundary maps $b_{ie,x}:W_{i,x}\to B_{e,x}$ on each incident edge. Let
\[
 Z_x=\{(w_i)_i: b_{ie,x}(w_i)=b_{je,x}(w_j)\text{ for every }e=ij\},
\quad u_i:Z_x\to W_{i,x}.
\]
The full coupling set is
\[
 \Coup_x=\{q\in TZ_x:Tu_i(q)=p_{i,x}\text{ for all }i\}. \tag{C}
\]
An admission policy can specify a smaller $\Adm_x\subseteq\Coup_x$. A coupling is retained as data, not selected by an independence convention. Under weak-pullback preservation of $T$, matching edge effects give nonempty $\Coup_x$ by the existing tree-admission theorem. Restricted admission needs its own existence evidence. Cycles can obstruct existence even with matching marginals. Neither acyclicity nor weak-pullback preservation is needed to transport an already supplied coupling.

Suppose component comparison maps $r_{i,x}:W_{i,x}\to W'_{i,hx}$ preserve the component effects and commute with the actual boundary maps through $k_{e,x}:B_{e,x}\to B'_{e,hx}$. They induce
\[
 R_x:Z_x\longrightarrow Z'_{hx},\qquad R_x((w_i)_i)=(r_{i,x}w_i)_i. \tag{R}
\]
The target carrier here is the full matching-tuple set. Bodies are supplied as maps $\phi_{i,x}:W_{i,x}\to Y_{i,x}$ to an output compatibility diagram. Their boundary equations give a joint body $\phi_x:Z_x\to Y_x=\lim_i Y_{i,x}$. Compatible component output maps $k_i$ give $k:Y_x\to Y'_{hx}$. Require the primitive body equations $\phi'_{i,hx}r_{i,x}=k_i\phi_{i,x}$; these do not follow from effect or boundary preservation. This includes the existing joint-occurrence construction $Y_i=SX_i$ when $S$ preserves the relevant tree limit. The full output projections, not just public observations, must jointly determine the body.

\begin{lemma}[Local transport, and the exact fixed-choice obstruction]\label{lem:transport}
For every $q\in\Coup_x$, $TR_x(q)\in\Coup'_{hx}$ and $\phi'_{hx}R_x=k\phi_x$. For specified $q_x,q'_{hx}$, there is a native comparison over the specified component maps if and only if
\[
 TR_x(q_x)=q'_{hx}. \tag{L}\label{eq:local}
\]
Its raw map is necessarily $R_x$. On all retained choices these transports preserve identities and composition, subtree restriction and retained reassociation.
\end{lemma}
\begin{proof}
Boundary commutation makes $R_x$ well typed. Since $u'_iR_x=r_{i,x}u_i$, functoriality gives
\[
 Tu'_i(TR_xq)=Tr_{i,x}(Tu_iq)=Tr_{i,x}(p_{i,x})=p'_{i,hx}.
\]
Projecting the proposed body square to each $Y'_i$ gives the component body square. Joint monicity of the limit projections proves it. Any map over the component maps has the same projections as $R_x$, hence equals $R_x$ by joint monicity of the target witness projections. Its native effect equation is therefore exactly \eqref{eq:local}; conversely that equation and the derived body square give the required native cell. Composition, restriction and reassociation are equality on all tuple coordinates; applying $T$ gives the corresponding effect identities. No inverse comparison or canonical joint effect was used.
\end{proof}

\begin{corollary}[Categorical form on retained choices]\label{cor:choice-category}
Fix the overlap shape. Let $\mathcal D_T$ be the category of local effect/boundary diagrams and their effect-preserving, boundary-commuting maps. Then $D\mapsto\Coup(D)$, $r\mapsto TR$, is a functor $\mathcal D_T\to\Set$. Its category of elements has objects $(D,q)$ and arrows $r:(D,q)\to(D',q')$ precisely when $TR(q)=q'$. Projection to $\mathcal D_T$ is a discrete opfibration. A restriction of the object sets satisfying \eqref{eq:admission} is a subfunctor and has the same property.
\end{corollary}
\begin{proof}
Lemma~\ref{lem:transport} gives well-defined maps and the functor laws. For each $r:D\to D'$ and source object $(D,q)$, exactly one arrow lies over $r$, namely the one with target $(D',TRq)$. This is the discrete-opfibration lifting property; its composite law is functoriality. Restriction works exactly when the image remains admitted.
\end{proof}

This statement concerns canonical tuple choices with specified component legs. It does not identify this category with the larger category of arbitrary native cones, nor does it assert an opfibration on independently chosen sections over changing input interfaces. The distinction is essential to the following descent condition.

\paragraph{The additional local hypothesis.} For the full admission \eqref{eq:local} is solved constructively by retaining $(x,q)$ and mapping it to $(hx,TR_xq)$. For restricted admission the exact forward condition is
\[
 q\in\Adm_x\ \Longrightarrow\ TR_xq\in\Adm'_{hx}. \tag{A}\label{eq:admission}
\]
Here admissions are sets with all native transport maps allowed by the displayed data. For a non-full category of admitted cones, preservation of its allowed arrows is additional; object closure alone does not imply it. Condition \eqref{eq:admission} is local at a declared rule, not a mission-success condition.

For a fixed chooser $x\mapsto q_x$, transporting objects is weaker than transporting a section. A target chooser exists on $h(A)$ exactly when
\[
 hx=hy\ \Longrightarrow\ TR_xq_x=TR_yq_y \tag{D}\label{eq:descent}
\]
in the same target fibre. Necessity is immediate. Sufficiency defines $q'_{hx}$ to be the common value, which is admitted by \eqref{eq:admission}; outside the image of $h$ a chooser must still be supplied if required. An independently fixed target chooser must satisfy \eqref{eq:local}, a stronger specification than merely existing. A retained common command or coupling parameter can keep choices distinct, but its availability to a scheduler or controller requires a separate information justification. The analyst's witness is not automatically an actor's observation.

\paragraph{Counterexample: equal marginals do not suffice.} Let $T=\mathcal D_{\rm fin}$, two component witnesses be fair bits, and the overlap boundary be a singleton. The diagonal coupling assigns probability $1/2$ to each of $(0,0),(1,1)$; the independent coupling assigns $1/4$ to each pair. Identity component comparisons force $R=\id_{\{0,1\}^2}$. Thus no native comparison over those components sends the diagonal choice to the independent choice. Both are admitted by the same tree. If two source states choose diagonal and antidiagonal laws and $h$ forgets that state distinction, \eqref{eq:descent} fails. This disproves an unrestricted rule-format theorem from marginal comparisons alone.

\paragraph{Noncanonical presentations.} If the target cone has extra witnesses not determined by its component projections, uniqueness of $R$ fails. The correct local obligation is an actual lift $a:Z\to Z'$ with $u'_ia=r_iu_i$, $Ta(q)=q'$ and the full body square. With a singleton source and a fair auxiliary target bit, component projections may all agree while no deterministic native lift exists. A span comparison would be different data. If the joint body carries a further audit bit forgotten by all component output projections, even effect equality and witness uniqueness need not give its body square. These are explicit limits of the iff in Lemma~\ref{lem:transport}; arbitrary cone data remain retainable through supplied full native lifts.

\section{A finite rule format with one shared world}

Fix a configuration interface $\Sigma$ containing one actual shared world and the local states of the active components. The rule grammar below generates finite transition derivations. It does not restrict execution length, and it does not make a whole method atomic. The paired source and target configurations have a pure map $h:\Sigma\to\Sigma'$. Each enabled rule and command has a step of form \eqref{eq:native} whose body returns the next configuration and retained event/content. Comparisons are required at every admitted input and command, including raw histories of zero probability when histories are part of the input.

The constructors are primitive steps, pure transformations and framing, cases whose selectors commute with the input comparison, local joins of Lemma~\ref{lem:transport}, and dependent sequential substitution. Case branches have a common post-interface and content grade with common comparison maps, or specified tagged coproducts and their induced maps. Source-enabled rule labels must map to enabled target labels. A scheduler selects an enabled labelled rule; its selected label is retained. For a chosen stochastic scheduler its own effect and comparison belong to the primitive data. For a demonic scheduler the assertion ranges over every admitted choice and transports that choice; it does not give the comparison permission to change an adversary's move after seeing its outcome.

\paragraph{Actual framing.} For example, $\Sigma=W\times L_1\times L_2$. A primitive of component 1 takes $(w,l_1)$ and returns $(w^+,l_1^+,q)$. Its frame returns $((w^+,l_1^+,l_2),q)$: component 2 retains its local state while receiving the same new shared world. If $h=h_W\times h_1\times h_2$ and the primitive square holds, this framed square follows by its three projections. More generally the same argument works for declared compatible-state limits and pure frame maps. Treating the frame as an identity on component 2's old shared world would be incorrect.

\paragraph{Joint bodies and physical interpretation.} Local joins use complete boundary outcomes and full joint body projections. Alternatively a particular substrate can supply a local effect assembler with a proved naturality/conservation law, which then becomes an explicit additional primitive hypothesis. Two summaries ``spend 7 from stock 10, leaving 3'' do not establish a physical joint payment of 14: equal post-state labels omit which debits or reservations occurred. An admitted boundary must retain the actual shared event identity, or a local assembler must account for both effects. The theorem does not infer this law from marginal equality, and a whole-system kernel cannot be relabelled a local assembler to discharge it. Ordinary interference between successive microsteps needs no simultaneous resource merger at all.

For clarity, dependent substitution of $P:A\to B$ and $Q:B\to C$ has witnesses
\[
 W^{PQ}_a=\sum_{w\in W^P_a}W^Q_{b(w)},\qquad
 p^{PQ}_a=p^P_a\mathbin{\gg=}\bigl(w\mapsto p^Q_{b(w)}\mathbin{\gg=}(z\mapsto\eta(w,z))\bigr), \tag{B}\label{eq:bind}
\]
where $b(w)$ is the actual next interface returned by $P$. Its content retains the ordered pair of contents, or their declared writer accumulation. In the latter case content comparisons must preserve the unit and accumulation; for a fixed monoid this means monoid homomorphisms, not arbitrary set maps. No continuation law is averaged over an erased joint input.

\begin{theorem}[Generated finite comparison]\label{thm:rules}
Supply primitive native comparisons, commuting pure/case/frame maps, and the boundary/body data of Lemma~\ref{lem:transport} at every local join. Use transported retained choices satisfying \eqref{eq:admission}; for fixed choosers require \eqref{eq:descent} and use its descended choice, or check \eqref{eq:local} against the given target. Then every finite expression of the rule grammar has a constructible native comparison. The comparisons respect identity, vertical composition and the retained reassociation/unit laws. Every finite scheduled execution has a comparison preserving its retained chronology and shared-state identity; the raw comparison maps commute with prefix truncation. Effect compatibility with truncation is a separate affine conclusion below. No whole-system joint kernel square is a premise.
\end{theorem}
\begin{proof}
Induct on rule formation. Primitive steps are the supplied data. Pure maps have singleton witnesses and Dirac/unit effects; their equations reduce to the pure squares. Frames follow by projections as above. A commuting case selector picks the same branch, so the branch map supplies the comparison. Local joins are Lemma~\ref{lem:transport} with the chosen admission condition.

For substitution let $r_a$ compare $P$, and $s_b$ compare the following $Q$ over the corresponding intermediate pure map $k$. Define the dependent-pair map
\[
 (w,z)\longmapsto (r_aw,s_{b(w)}z).
\]
The first body square identifies the target continuation input as $k(b(w))$, so this is well typed. Naturality of bind and $Ts_b(p_b^Q)=p_{kb}^{Q'}$ give the pushforward of \eqref{eq:bind} as
\[
 p^P_a\mathbin{\gg=}\bigl(w\mapsto p^{Q'}_{kb(w)}\mathbin{\gg=}(z'\mapsto\eta(r_aw,z'))\bigr).
\]
The integrand depends on $w$ only through $r_aw$: the target intermediate body is $b'(r_aw)=kb(w)$. Using $Tr_a(p^P_a)=p^{P'}_{ha}$ therefore yields exactly the target dependent-pair effect. The body square follows from the two body squares and naturality of ordered writer accumulation.

All witness maps act on actual coordinates, so identity and vertical composition follow coordinatewise. Monad laws and Cartesian reassociation give the unit/associativity comparisons and their coherence; no exchange of independent effects is used. Thus no commutativity assumption on $T$ is necessary. A finite execution is a repeated dependent substitution of the selected rules. Its comparison follows by a second induction on prefix length, using the actual post-configuration at each step. The maps commute with dropping the last raw step. This proves the prefix claim, including every admitted scheduler word.
\end{proof}

The hypothesis is a formation condition on generators and local constructors, not an assumption that composites already have comparisons. The proof gives the witness map and its certificate. It has substance only when primitive boundaries are genuinely local and rule formation does not hide future interference or the desired global guarantee. If a syntax is later quotiented by further equations, its comparison must respect those equations; arbitrary method permutations and trace equations are not among the proved equations.

The generated carrier in \eqref{eq:bind} retains its intermediate shared state. It is not silently replaced by a canonical tuple of component composites matching only their final boundaries: that larger tuple can combine histories with incompatible intermediate shared states. The existing directed sequential interchanger is not inverted here. Flattening to component histories requires all intermediate boundary traces or a separate exactness theorem.

\paragraph{A nonserializable instance.} Two increments of shared $x$ have separate read and write microsteps. Starting at $0$, the scheduler word $R_1R_2W_1W_2$ stores two local reads of $0$ and ends at $1$. A whole-method serial word $R_1W_1R_2W_2$ ends at $2$. Both remain generated histories. A primitive witness may additionally retain an auxiliary bit erased by a directed comparison; the state/event/body squares preserve the actual read/write schedule and final state. This is interference at a declared microstep boundary, not a proof that a nonserializable method admits serial replay. The executable diagnostic checks all six read-before-write shuffles.

\section{Raw coverage and adversarial quantifiers}

Effect preservation does not imply raw surjectivity. Let both source witness sets be $\{0,1\}$, sharing the identity boundary. Collapse the boundary to a singleton while leaving both component witness maps identity. The source joint carrier is the diagonal and the target carrier all four pairs. Each component map is surjective, but the joint map is not. Choosing the target diagonal effect gives a native comparison while leaving two target raw witnesses uncovered.

\begin{proposition}[A local sufficient coverage condition]
Root the overlap tree. Suppose the root comparison is surjective and each child $i$, at its parent edge $e$, has a surjective relative map
\[
 (r_i,b_{ie}):W_i\longrightarrow W'_i\times_{B'_e}B_e.
\]
Then the induced joint map $R:Z\to Z'$ is surjective.
\end{proposition}
\begin{proof}
Lift a given target root witness. Its source boundary fixes the parent-edge value needed by each child. Since the target tuple matches and boundary squares commute, the target child witness and that source value form an element of the displayed pullback. Lift it and recurse down the finite tree. The resulting tuple matches every edge and maps to the given target tuple.
\end{proof}

Restricted raw admission constraints require compatible lifts too. Along executions, local coverage is needed at each already chosen source preimage of the current target state. Together with initial-state coverage and matching rule availability it lifts finite target paths; coherent successive choices give infinite raw lifting under dependent choice. Surjectivity at every unrelated finite horizon alone is insufficient: the inverse sequence $P_n=\{m\in\mathbb N:m\geq n\}$ with inclusion bonds maps surjectively to the constant singleton sequence but has empty inverse limit. For transferring a target universal property back along a total source-to-target comparison, target coverage is not needed; every source path already has its image. Coverage is required when the claimed transfer quantifies over paths outside that image or uses a comparison span whose source projection must cover all source cases.

\section{Test against the original relative presentation}

For the effect-fibre context category over $A$, an object is a family $(W_a,p_a)$ and a native map is the actual family $r_a$ with $Tr_a p_a=p'_a$. Its extension is $e_AP=\sum_aW_a$. For one-hole writer content with grade $Q$, take $d_BQ=B\times Q$ and actual content maps $v:Q\to Q'$. A mixed operation is exactly
\[
 (P,Q,f:e_AP\to d_BQ)\in H(A,B)=(e_A\downarrow d_B).
\]
Thus \eqref{eq:native} supplies the original native comma cell, including its context and content maps. For the effect class of the existing retained-fibre package this is its terminal-base/base-preserving specialization, where the comma presentation is literal; the arbitrary-monad finite sequential generalization is justified by the direct proof here. No synchronized-tensor result is imported under the weaker assumptions. Over the full bundle category the existing construction uses normalization; this note asserts no unrestricted raw-comma equivalence there.

The mixed crossing retains a writer value while reindexing the next context at the actual output: $(w,z,q)$ and $(w,q,z)$ are the same retained data under the specified tuple bijection. The next context depends on the returned interface, which includes the current shared world, rather than on an unrecorded effect. Writer $Q$ alone is not an operative store: if a pending effect changes a continuation, it must also be represented in the returned input/state $B$, or another interaction law must be supplied. The ordinary writer crossing reindexes along the projection dropping $Q$. Context accumulation is \eqref{eq:bind}; writer accumulation is ordered Cartesian pairing (or a specified monoid such as event words, with unit/accumulation-preserving comparison maps). The unit and accumulation squares are the tuple/monad calculations used in Theorem~\ref{thm:rules}. This tests the construction at the original native presentation, not only at its observable kernel. There is one continuation hole; arbitrary multi-occurrence shared continuation is not inferred from an independent occurrence tensor.

\paragraph{Why the general comma shape is insufficient.} Take context grades to be subsets of $\{a,b\}$, ordered by inclusion, with union tensor and trivial action on $\Set$; take trivial content. Over a singleton interface all carrier extensions and bodies are identical. Grades $\{a\}$ and $\{b\}$ both map to $\{a,b\}$, but there is no native grade map $\{a\}\to\{b\}$. Carrier/effect/body equality cannot manufacture it. Consequently the effect-fibre criterion cannot be exported to an arbitrary original context/content doctrine by forgetting grades.

\paragraph{The remaining categorical lifting problem.} For an independently supplied doctrine, one needs a local rule interpretation on its actual category of admitted native cones: primitive comparison and boundary data must produce a directed joint context/content cell, with a choice of target cone in the admitted class, compatibly with pure reindexing and the specified frame, case and dependent-substitution operations. For fixed cone sections the lift must also descend. Identities, composites and the syntax equations must be preserved by these lifts. These are local native lifting and coherence obligations, not a whole-system kernel square. In the canonical effect-fibre fragment the tuple construction proves them. In arbitrary grades they require a genuine grade-lifting construction; the subset-grade example shows why carrier calculations alone cannot discharge it. Section~\ref{sec:grade-lifting} supplies such a theorem for a parameterized class using independent coefficient limits, without assuming the mixed lifts. A universal opcartesian lift is stronger than the comparison construction requires and is not claimed. There is no established uniquely weakest hypothesis over all doctrines.

Existing hierarchical cone substitution consumes supplied native legs; it does not choose a preserved joint law. Existing shared-witness products provide conjunction with the same shared state under their exactness hypotheses; they do not enumerate scheduler words or select the physical coupling. Neither mechanism removes \eqref{eq:descent}, changes a directed cell into an isomorphism, or makes independently closed copies of shared state identical.

\section{Test against information and progress}

For $T=\mathcal D_{\rm fin}$, let a generated rule have effect $q_\sigma$ and full body $A_\sigma:Z_\sigma\to\Sigma\times Q$. Its actual kernel is $K(\sigma)=T A_\sigma(q_\sigma)$. Theorem~\ref{thm:rules} gives $R_\sigma$, its effect equation and $A'_{h\sigma}R_\sigma=(h\times v)A_\sigma$. Hence
\[
 T(h\times v)K(\sigma)
 =T A'_{h\sigma}(TR_\sigma q_\sigma)
 =T A'_{h\sigma}(q'_{h\sigma})
 =K'(h\sigma). \tag{J}\label{eq:global}
\]
This derives the actual joint next-state/event equation required by condition 3 of the existing Information-and-Progress theorem for systems and partners generated by the same rule format. A compatible partner contributes its own local rules, boundary/choice data and information interface; when comparison is identity on its private part the same induction yields the actual shared-world product map. No assumption of an arbitrary enlarged-system kernel equation is necessary. A partner with a different independently specified joint chooser or unmodelled effect can violate the local hypotheses, and is not covered merely because its marginals agree.

The rest of that theorem is not absorbed into \eqref{eq:global}. Equal event and task menus require guard reflection as well as forward preservation; full rule/firing identities must be retained. Source-local information must factor through target-local information, at actual product histories, for source policies to induce admissible target policies. This is stronger than the common-observation condition for one supplied controller. Initial effects must push forward. Goal reflection and receipts remain body obligations. Weak/strong fairness uses the reflected enabledness and firing data; physical-time progress, erased-step productivity, non-Zeno behavior and deadlines need their own evidence. The construction supplies the missing generated plant comparison and compatible raw maps, not a Byzantine completion theorem.

\section{Guarded feedback: what is and is not independent}

\begin{corollary}[Finite guarded preservation]\label{cor:feedback}
Suppose the rules of Theorem~\ref{thm:rules} are available at every admitted state/history and command. A strictly delayed controller using a preserved observation prefix induces native comparisons at every finite closed prefix. For affine $T$ the prefix effects are compatible with truncation. This needs no independent interaction axiom beyond the local conditional native equations and the controller observation square.
\end{corollary}
\begin{proof}
Let $P_0=1$, and let $P_{n+1}=\sum_{p\in P_n}W_{s(p),a(p)}$, with the effect obtained by bind. At matched prefixes the controller sees the same observation and chooses the same command. Define
\[
 R_{n+1}(p,w)=(R_np,r_{s(p),a(p)}w).
\]
The body equation makes the target input well typed. Bind naturality gives $TR_n\theta_n=\theta'_n$ by induction, and coordinate deletion gives $p'_nR_{n+1}=R_np_n$. If $T1\cong1$, discarding the last conditional effect yields the unique terminal effect, so $Tp_n\theta_{n+1}=\theta_n$. Without affinity, same-horizon comparison remains valid but this discard law requires separate evidence. A delayed wire is included by keeping its old value in state and storing the newly returned value for the next step.
\end{proof}

This is an application of the existing retained-fibre guarded-prefix theorem. A whole-tape comparison preserving the law for each fixed open input is insufficient: the native-first-results manuscript gives a sectionwise measure-preserving comparison whose exceptional null diagonal is selected with certainty by feedback. Our all-state/all-command premise excludes precisely that failure. Stochastic controllers require their own retained innovations and conditional native maps; existential couplings without the necessary indexed choice or measurable selection do not suffice in richer measurable doctrines.

\paragraph{An independent infinite-effect issue.} Compatible raw prefix maps define a map of infinite raw histories coordinatewise. They need not give an effect on those histories. For $T=\mathcal D_{\rm fin}$, let $\theta_n$ be the uniform law on $\{0,1\}^n$. The family is truncation-compatible. An infinite-path law in $\mathcal D_{\rm fin}(\{0,1\}^{\mathbb N})$ would have some finite support size $k$, so every prefix marginal would have support at most $k$, contradicting $2^n>k$. Thus the same affine monad satisfies the finite theorem and fails infinite-effect realization.

If the intended feedback object consumes a completed infinite effect, one needs an additional completion result: the particular compatible family must lie in the image of
\[
 T(\lim_nP_n)\longrightarrow\lim_n TP_n,
\]
with suitably compatible choices or a uniqueness property for transport. Preservation of the relevant inverse limits is a strong sufficient condition. The existing continuous display-completion theorem supplies a separate conditional route in its actual presentation category. This need not be added for a task using only all-raw path simulation and finite effects. Instantaneous feedback, trace axioms, fairness and physical-time divergence are further distinct questions, not consequences of guardedness.

\section{Native grade lifting from independent coefficient limits}\label{sec:grade-lifting}

The preceding effect-only theorem extends to a class with nonfaithful context and content extensions. The joint native arrows are now constructed, rather than postulated as lifts of their carrier maps. The class includes an original augmented-rig doctrine exactly. It remains a parameterized family theorem: arbitrary contextad/contentad pairs are not being classified.

\subsection{Independent one-sided inputs and the original mixed presentation}

Fix a monad $T$ on $\Set$ and locally small monoidal categories $(\mathcal V,\otimes,I)$ and $(\mathcal D,\odot,J)$. Neither symmetry nor a relationship between the two monoidal structures is required. Choose limits for the finite coefficient diagrams used below. Finite completeness of $\mathcal V,\mathcal D$ is a convenient sufficient envelope; only the actual boundary diagrams, their comparison images and their retained iterated decompositions need limits. No extension functor is assumed faithful, full, cartesian-lifting, or opcartesian-lifting. In particular no limit or desired comparison in a mixed hom category is an input.

Write $\operatorname{Fam}(\mathcal D)$ for the category whose objects are pairs $Q=(O,(D_o)_{o\in O})$ and whose arrows are
\[
 (v,\beta):Q\longrightarrow Q',\qquad
 v:O\to O',\quad \beta_o:D_o\to D'_{v(o)}.
\]
Composition is $(v',\beta')(v,\beta)=(v'v,(\beta'_{v(o)}\beta_o)_o)$. The carrier functor forgets the actual coefficient maps and is generally nonfaithful. Cartesian products of carriers and the given tensor define
\[
 Q\boxtimes R=(O\times O_R,(D_o\odot E_r)_{(o,r)}),\qquad
 \mathbf 1_{\mathcal D}=(1,J).
\tag{GW}\label{eq:graded-writer}
\]
This is a monoidal category with the actual tensor associators and unitors on coefficients, not merely carrier bijections.

The context fibre $\mathcal M_A$ consists of families
\[
 P_a=(W_a,p_a,(C_{a,w})_{w\in W_a}),\qquad p_a\in TW_a,\quad C_{a,w}\in\mathcal V.
\]
A native arrow over a pure input map $h:A\to A'$ consists of
\[
 r_a:W_a\to W'_{ha},\qquad
 \alpha_{a,w}:C_{a,w}\to C'_{ha,r_aw},\qquad
 Tr_a(p_a)=p'_{ha}. \tag{G}\label{eq:graded-context-map}
\]
All of these arrows are retained. In particular two choices of $\alpha$ over the same $r$ remain different native arrows. The extension is $e_AP=\sum_aW_a$, and pure reindexing is pointwise substitution. The context unit has $W_a=1$, effect $\eta(*)$ and coefficient $I$. If a continuation $R$ is indexed by $(a,w)\in e_AP$, accumulation has
\[
 W^{P\star R}_a=\sum_{w\in W_a}W^R_{a,w},\qquad
 C^{P\star R}_{a,(w,z)}=C_{a,w}\otimes C^R_{a,w,z},
\tag{GA}\label{eq:graded-accumulation}
\]
and the retained bind effect of \eqref{eq:bind}. On a native dependent pair its coefficient map is $\alpha_{a,w}\otimes\gamma_{a,w,z}$, with the continuation map taken at the actual retained input. Composition uses the actual coefficient composites, as well as the witness-map composite.

For content take $\mathcal N_B=\operatorname{Fam}(\mathcal D)$ and $d_BQ=B\times O$. Pure output transport changes $B$ and leaves $Q$ fixed. Content unit and accumulation are \eqref{eq:graded-writer}. Thus $p:\Set\times\operatorname{Fam}(\mathcal D)\to\Set$ and $d(B,Q)=B\times|Q|$ give the independent writer content presentation, including all noninvertible coefficient arrows.

\begin{proposition}[Independent presentations and derived writer crossing]\label{prop:graded-presentation}
These data define a contextad and a contentad and a native mixed interaction on their original literal comma family
\[
 H(A,B)=(e_A\downarrow d_B).
\]
An object is $(P,Q,f:e_AP\to B\times O)$; a changing-endpoint native cell is an actual context map $(r,\alpha)$, an actual content map $(v,\beta)$, and the pure equation $f'r=(k\times v)f$, where $r(a,w)=(h(a),r_a(w))$ is the total extension map. The crossing and composition retain $\alpha,\beta$, including when their carrier maps are identities.
\end{proposition}
\begin{proof}
Context reindexing is strictly pointwise. For accumulated contexts, the native effect equation is bind naturality as in Theorem~\ref{thm:rules}. Coefficient functoriality is
\[
 (\alpha'\otimes\gamma')(\alpha\otimes\gamma)
 =(\alpha'\alpha)\otimes(\gamma'\gamma).
\]
The continuation indices on both sides agree by the retained witness map. Context associators retain the ordered triple $(w,z,t)$ and use the associator of $\mathcal V$; their effect equation is monad associativity. Unit maps use the monad units and the unitors of $\mathcal V$. Naturality, pentagon and triangle are respectively tensor naturality and the actual monoidal axioms. These verify the one-sided context structure on its native category, without quotienting coefficient arrows. The content structure follows from the monoidal category \eqref{eq:graded-writer}; its extension accumulation is Cartesian reassociation and its coefficient accumulation is the given tensor in $\mathcal D$.

Construct the crossing before constructing any mixed join. A content-then-context representative has a writer body $a\mapsto(b(a),o(a))$, grade $Q$, and a following context $R$ over $B$ with a pure body $g:e_BR\to C$. Cross it to the context over $A$ with fibre $W^R_{b(a)}$, effect $p^R_{b(a)}$ and coefficients $C^R_{b(a),w}$; retain the content grade $Q$ and use the body
\[
 (a,w)\longmapsto\bigl(g(b(a),w),o(a)\bigr). \tag{GC}\label{eq:graded-crossing}
\]
On arrows this uses the supplied context coefficient map at $b(a)$ and the supplied content coefficient map at $o(a)$. Their pure body equation supplies exactly the reindexing equality. Consequently native naturality holds on the actual pair of grade arrows, not just on extensions. A pure middle change either postcomposes $b$ first or pulls $R$ back first; both select the same fibre, effect and coefficient. This proves pure-middle balance and its identity/composition equations.

The four crossing compatibilities are explicit. Crossing the context unit retains $Q$ and uses the context singleton/unit coefficient. Crossing the writer unit removes its singleton content coordinate. Crossing two writer stages retains their ordered output pair and the same $\mathcal D$ tensor; crossing two context stages retains their dependent witness pair, the same bind effect and the same $\mathcal V$ tensor. The comparisons are the one-sided unitors, associators and the indicated pure reindexing identifications. For further composites every route retains the same addressed witness and writer coordinates. On effects its equations are the monad laws; on the two independent coefficient categories they are their monoidal naturality, pentagon and triangle laws. Thus the mixed compatibility diagrams commute on actual grades. No tensor is moved across the other coefficient category or across an unrelated effect. In particular no commutativity of $T$ and no tensor-limit preservation are used.

The mixed multiplication therefore has the dependent witnesses and effect of \eqref{eq:bind}, context coefficient $C_w\otimes C^R_{b(w),z}$, content grade $Q\boxtimes Q_R$, and the full returned body. The original retained comma factorization through $e_AP$, together with the pure-middle balance just checked, supplies the original relative presentation. There is no new span quotient or carrier-only replacement of $H$.
\end{proof}

Only one continuation hole is asserted. An operative pending effect must remain in the returned state/input, as before. The graded writer crossing does not permit a later context to inspect forgotten writer data. General polynomial occurrences or other independent content mechanisms need their own interaction proof; the present proposition is an independently constructed writer class.

\subsection{Coefficient limits construct the joint native cells}

Let $\mathcal I$ be a finite nonempty boundary diagram, for example vertices and edge objects of a finite overlap tree. Supply a diagram of primitive mixed objects $(P_i,Q_i,f_i)$ with inputs $A_i$, outputs $B_i$, and actual native boundary arrows. Thus a boundary arrow supplies its pure input/output maps, its witness map, its $\mathcal V$ coefficient maps, its $\mathcal D$ content maps, and its body equation. Supply a comparison of two such diagrams: actual primitive native comparisons commuting with those boundary arrows in both grade categories. Merely commuting carrier squares are insufficient.

Form $A=\lim_i A_i$ and $B=\lim_i B_i$. At $a=(a_i)\in A$, the boundary witness maps define a diagram of sets and its matching-tuple carrier
\[
 Z_a=\lim_i W^i_{a_i},\qquad u_i:Z_a\to W^i_{a_i}.
\]
Retain a coupling $q_a\in TZ_a$ with $Tu_i(q_a)=p^i_{a_i}$. Coupling existence and restricted admission are exactly the earlier separate obligations. For each $z=(w_i)\in Z_a$, the actual boundary coefficient maps define a diagram $\mathcal C_{a,z}:\mathcal I\to\mathcal V$. Define
\[
 P^{\mathcal I}_a=(Z_a,q_a,(\lim_{i\in\mathcal I}\mathcal C_{a,z})_{z\in Z_a}). \tag{GJ}\label{eq:joint-context-grade}
\]
Its native projections use $u_i$ and the coefficient-limit projections. They are derived from primitive grade data and the coupling marginal equations.

There is a direct limit construction in $\operatorname{Fam}(\mathcal D)$: for a diagram $Q_i=(O_i,D^i)$ put
\[
 O=\lim_iO_i,\qquad
 Q^{\mathcal I}=\bigl(O,(\lim_i D^i_{o_i})_{o=(o_i)\in O}\bigr). \tag{GL}\label{eq:family-limit}
\]
Here the coefficient arrows in each limit are those of the supplied native content diagram. To verify the native universal property, a cone with domain $(L,E)$ first gives a unique carrier map $L\to O$. At each $l\in L$, its actual coefficient cone gives a unique arrow $E_l\to\lim_iD^i_{o_i(l)}$. These maps supply a unique family arrow. This proves \eqref{eq:family-limit} is a limit on the actual native category, not merely a carrier limit.

For $f_i(a_i,w_i)=(b_i,o_i)$, the boundary body equations imply that $(b_i)$ and $(o_i)$ are compatible tuples. Set
\[
 f^{\mathcal I}(a,z)=\bigl((b_i)_i,(o_i)_i\bigr)\in B\times|Q^{\mathcal I}|. \tag{GB}\label{eq:joint-body}
\]
This uses the full output limit: $B\times O\cong\lim_i(B_i\times O_i)$. For a common shared output $B_0$, a nonempty connected constant-$B_0$ diagram gives $B\cong B_0$; a disconnected family instead gives distinct output copies unless an explicit shared boundary is imposed. There is no identification of those copies by convention.

\begin{theorem}[Native grade lifting in the original presentation]\label{thm:grade-lifting}
Under the independent inputs of Proposition~\ref{prop:graded-presentation} and the selected coefficient limits above, a comparison of primitive boundary diagrams induces a directed joint native context/content cell between \eqref{eq:joint-context-grade}--\eqref{eq:joint-body} whenever its retained effects satisfy $TR_a(q_a)=q'_{ha}$. Here $h:A\to A'$ and $R_a:Z_a\to Z'_{ha}$ are induced by the primitive input and witness maps. No joint grade arrow, joint mixed cell, or whole-system comparison is an input. On retained admitted choices the construction is functorial. For fixed input-indexed choices, admission closure and descent remain required exactly as in \eqref{eq:admission}--\eqref{eq:descent}.
\end{theorem}
\begin{proof}
At each $(a,z)$, primitive coefficient comparisons form a natural transformation
\[
 \alpha_{a,z}:\mathcal C_{a,z}\Longrightarrow\mathcal C'_{ha,R_az}.
\]
Its naturality is precisely the actual coefficient part of boundary commutation. The source limit projections followed by its components form a cone over the target coefficient diagram. Target limit universality constructs the unique arrow
\[
 \overline\alpha_{a,z}:\lim_i\mathcal C_{a,z}\longrightarrow\lim_i\mathcal C'_{ha,R_az}.
\tag{LG}\label{eq:coefficient-lift}
\]
Thus $(R_a,\overline\alpha_{a,z})$ is an actual native context arrow: its effect equation is the stated retained-choice equality, and its coefficient maps were constructed independently of that equation. The content comparison is the actual map between the limits \eqref{eq:family-limit} induced by the primitive content arrows; denote it $(v,\overline\beta)$. The primitive body equations, tested against every $B'_i$ and $O'_i$ projection, give
\[
 f'^{\mathcal I}R=(k\times v)f^{\mathcal I}.
\]
This is the full mixed body equation. Together with the two constructed grade arrows it supplies the desired cell in the original $H$.

For an identity comparison the identity on each coefficient limit has the required projections, so uniqueness makes it the constructed map. For successive comparisons, the composite of the two constructed maps has the projected composites of the primitive maps; uniqueness identifies it with the map constructed from their composite. The same proof works in $\operatorname{Fam}(\mathcal D)$, and the witness/effect laws are ordinary functoriality. This proves identity and composition on actual native arrows, even when extension forgets distinct arrows. No faithfulness of $e$ or $d$ is used.
\end{proof}

Coefficient-limit universality does not make the chosen effect-bearing joint a categorical product in the entire mixed $H$: retained couplings can differ, and a universal mediating effect map need not exist.

For canonical raw carriers, the condition $TR_aq_a=q'_{ha}$ is still necessary for a comparison over the given primitive legs. In this coefficient-limit presentation the coefficient lift is automatically supplied by \eqref{eq:coefficient-lift}, and is unique over those legs. If one independently specifies joint grades that are not the constructed coefficient limits, the theorem does not manufacture maps between them. Their different admission/grade choice is additional data, not a consequence of equal carriers.

Actual primitive coefficient naturality cannot be dropped. With all raw carriers singleton, take source and target coefficient cospans $k\xrightarrow{\id}k\xleftarrow{\id}k$ in vector spaces. Choose the left component comparison $\id$ and the right comparison $2\id$. No boundary comparison makes both squares commute, and no map between their diagonal limits projects to those two maps. All carrier and pure body equations are nevertheless trivial. The failed obligation is visible in a local coefficient square, before any whole-system comparison.

\subsection{Coherence of the generated constructors}

\begin{theorem}[Native rule closure and coherence]\label{thm:graded-rules}
For full coupling admission, the constructed comparisons are closed under pure input reindexing, compatible shared-state framing, typed cases, retained local joining, and dependent substitution. A restricted policy must include the particular choices used by these constructors and their comparison transports; this is a local formation condition, not inferred from marginal correctness. They preserve identity and vertical composition and commute with the specified unit, reassociation and pure-middle balance comparisons. Joining and substitution admit the canonical forward comparison below; no inverse or arbitrary distributivity law is asserted.
\end{theorem}
\begin{proof}
\emph{Reindexing.} At $a_0\in A_0$, reindexing along $s:A_0\to A$ selects the exact tuple, effect and coefficient diagram at $s(a_0)$. Thus performing the joint construction before or after this reindexing gives the same local data, up to the chosen limit comparison. Its native comparison is determined by the same projections. Identity/composite reindexing follows by substitution and uniqueness. This statement compares component restrictions along the projections of the actual $s$ and retains the common parameter $A_0$. For a disconnected shape, the literal joint of the restricted input carriers has base $\lim_i\operatorname{const}(A_0)$; restrict its raw/grade construction along the diagonal $A_0\to\lim_i\operatorname{const}(A_0)$ before assigning effects. Only the existing choice $q_{s(a_0)}$ is used; no off-diagonal coupling choices are required. No extra copies of the input parameter are identified silently. No new independently chosen coupling section is created.

\emph{Framing.} A pure frame carries additional private coordinates while the declared shared world is updated once. It changes neither the witness effect nor its coefficient. Apply the original pure input/output actions to the full body, retaining the coefficient arrows. Projection to each world/private/content coordinate proves the framed body equation. Successive frames associate by pure Cartesian reassociation; their coefficient maps are unchanged and therefore satisfy the same equation. Compatible-state pullback frames use their actual pullback projections. No duplicated old copy of the shared world is held fixed.

\emph{Cases.} A selector commuting with the input comparison chooses corresponding branches. Context grades, effects and coefficient maps are selected pointwise. Branches either share a common content grade and output interface with common comparison maps, or are embedded in specified tagged coproducts. In $\operatorname{Fam}(\mathcal D)$ the latter coproduct has carrier the disjoint union and the original coefficient on each tag; its native universal property follows branchwise. Preservation of identity/composition and nested-case reassociation is tested on each retained tag. These facts do not assert that a case selector may be moved across a join when it depends on forgotten component data.

\emph{Dependent substitution.} For two steps the context map on witnesses is $(w,z)\mapsto(rw,s_wz)$ and on coefficients is $\alpha_w\otimes\gamma_{w,z}$. The native first-body square identifies the continuation input. Bind naturality proves the effect equation, and the actual tensor bifunctor proves the coefficient identity/composition law. Content maps are $\beta_o\odot\delta_p$. Their source and target grades are the actual ordered content products. Hence the mixed body square follows from the two constituent body squares. Triple substitution uses the same retained witness triple, monad associativity, the associator of $\mathcal V$ on context coefficients and that of $\mathcal D$ on content coefficients. Their naturality, pentagon and unit triangles prove the native reassociation laws. No step identifies a coefficient map from its carrier image.

\emph{Hierarchical joining.} Decompose a finite overlap tree into connected subtrees, retaining the shared separator objects. Fix one coupling on the full matching-tuple carrier and use its projections for every subtree. Both nested raw matching constructions are canonically the full matching-tuple carrier; effects are its retained marginals. At a tuple, a cone to the full coefficient diagram is equivalent to cones to the subtree diagrams agreeing on the separator. Therefore the iterated coefficient limit and the flat coefficient limit represent the same cone functor. The unique comparison preserving all primitive projections is a native isomorphism. This argument applies separately in $\mathcal V$ and in $\operatorname{Fam}(\mathcal D)$. For any two reassociation routes, projecting to every primitive factor gives the same actual coefficient morphism. Limit uniqueness proves equality, including the pentagon and naturality under primitive comparisons. A one-vertex join has its identity limit comparison. The argument does not equate independently sampled intermediate couplings or discard intermediate boundary witnesses.

Finally induct on the finite constructor expression. Each case constructs the next actual pair of grade arrows and verifies its full body equation. The constructor identity/composition laws just proved make this induction functorial, and their displayed naturality makes it compatible with the admitted reassociation and unit laws. Any further quotient of syntax still owes preservation of its additional equations.
\end{proof}

\paragraph{Directed join/substitution comparison.} Consider two chronological layers of boundary-compatible primitive operations, with the second conditioned on the actual returned component input. First join each layer, then substitute. Its dependent history maps to the tuple of component dependent histories. At one such history its context coefficient map is induced by the cone
\[
 (\lim_i C_i)\otimes(\lim_i E_i)
 \xrightarrow{\pi_i\otimes\rho_i} C_i\otimes E_i.
\]
Thus it has the directed map to $\lim_i(C_i\otimes E_i)$. The content map is similarly the canonical map from the tensor of its two actual family limits to the limit of the component content tensors. No tensor-limit preservation hypothesis is needed to construct either map. Conditional bind and native projection equations show that the transported retained effect has the required component-composite marginals; use that transported choice at the target when admitted (automatic for full admission), or check equality against a separately fixed admitted choice. The full body equation follows componentwise. Consequently this is a derived native cell
\[
 \mathsf J(P_i)\,;\,\mathsf J(R_i)\longrightarrow\mathsf J(P_i;R_i), \tag{JI}\label{eq:graded-interchanger}
\]
with the stated retained raw carriers and target choice.

Its naturality is checked after every coefficient/content projection: both routes are the same tensor composite of the primitive maps and projections. Multiple chronological layers reduce in this way to the same ordered tensor of projections, with monad associativity on effects. Monoidal coherence and limit uniqueness prove the reassociation/composition diagrams. The unit comparison is the cone induced by the monoidal unit maps. For a nonempty connected constant-unit diagram its limit is canonically the same unit; for a disconnected shape only the diagonal map $I\to\lim_i I$ (and similarly $J$) is automatic. A nullary join, if included, uses terminal coefficient objects, not automatically the chronological units. These distinctions prevent an unsupported strong or invertible interchange claim.

Only the actual generated dependent carrier is used on the source of \eqref{eq:graded-interchanger}. If a target forgets intermediate shared-boundary agreement, its raw carrier can be larger. If its coefficient tensor fails the corresponding exactness, its grade map can also fail to be invertible even when the raw carriers agree. The example below demonstrates the latter failure.

\subsection{Recovery and a materially different original native doctrine}

Taking $\mathcal V=\mathcal D=\mathbf 1$ removes exactly the coefficient objects and arrows. Proposition~\ref{prop:graded-presentation} becomes the earlier monadic effect-fibre/writer presentation; \eqref{eq:joint-context-grade} retains all couplings; Theorems~\ref{thm:grade-lifting} and \ref{thm:graded-rules} give the earlier comparison and rule-format theorems with ordered content pairs. A further fixed writer-monoid accumulation is the additional unit/accumulation-preserving interpretation already specified above. No restriction to independent couplings has entered this recovery.

For a genuinely different native-grade instance, let $k=\mathbb Q$ and
\[
 \mathcal V=\operatorname{AugVect}_k=\operatorname{FinVect}_k/k.
\]
An object is $(C,\epsilon:C\to k)$, a map $a:(C,\epsilon)\to(C',\epsilon')$ is an actual linear map satisfying $\epsilon'a=\epsilon$, and the monoidal tensor is $(C\otimes C',\epsilon\otimes\epsilon')$, using $k\otimes k\cong k$. The monoidal unit is $(k,\id)$. This slice has finite limits. In particular its product is the actual linear pullback
\[
 (C,\epsilon)\times_k(E,\delta)
 =\bigl(\ker[\epsilon,-\delta],\ (c,e)\mapsto\epsilon c=\delta e\bigr). \tag{AV}\label{eq:augmented-pullback}
\]
Use limits in the slice: a disconnected diagram does not have the unqualified underlying vector-space product. Coefficient boundary pullbacks are likewise kernels of differences of the actual boundary linear maps, with their induced augmentation. A primitive comparison $(a,b)$ induces $(c,e)\mapsto(ac,be)$ on this kernel because its boundary square commutes. This explicitly constructs the joint grade map.

Take $T$ to be the constant-singleton monad and restrict raw witness carriers to finite sets. Its effects impose no additional data or equations. Then the context fibre above recovers the original augmented-rig family over $\operatorname{FinVect}_k$, not merely an analogous carrier model. Indeed an original augmented family has
\[
 (W,(C_w)_w,\epsilon:\textstyle\bigoplus_w C_w\to k).
\]
Because finite direct sums are coproducts, $\epsilon$ is uniquely its family of restrictions $\epsilon_w:C_w\to k$. Similarly an original native aggregation map
\[
 A_{w'}:\bigoplus_{r(w)=w'}C_w\longrightarrow C'_{w'}
\]
is uniquely its restrictions $\alpha_w:C_w\to C'_{r(w)}$. The global augmentation equation is equivalent to $\epsilon'_{r(w)}\alpha_w=\epsilon_w$ for every $w$. This gives an equivalence with $\operatorname{Fam}_{\rm fin}(\operatorname{AugVect}_k)$ retaining every native arrow. Under it, original dependent accumulation restricts on summand $(w,z)$ to the same tensor and augmentation $\epsilon_w\otimes\delta_{w,z}$ as \eqref{eq:graded-accumulation}; the aggregation maps restrict to $\alpha_w\otimes\gamma_{w,z}$. Reindexing, units and associators agree. Thus the equivalence respects the one-sided context operations, not just their object lists. It relies on the coproduct property and is not asserted for arbitrary rig addition.

Content can independently take $\mathcal D=\mathbf1$ or $\operatorname{AugVect}_k$. In the latter case its native arrows also carry linear maps forgotten by extension. This writer content has occurrence projection $B\times|Q|\to B$ for the identity occurrence monad; reindexing at a pure map changes only the base, so its display and base-change squares are the actual Cartesian squares. It is consequently compatible with the original one-hole occurrence presentation. No arbitrary new multi-occurrence content is inferred.

\paragraph{Nonfaithfulness and a strict directed interchanger.} Let $U=(k^2,\epsilon(x,y)=x)$. Both $\id_U$ and $s=\operatorname{diag}(1,0)$ preserve augmentation and have the same singleton raw witness map. They are distinct native arrows. The joint grade $P=U\times_kU$ is $k^3$ with coordinates $(a,b,c)\mapsto((a,b),(a,c))$. Applying $\id_U$ to the first leg and $s$ to the second constructs the directed joint map $(a,b,c)\mapsto(a,b,0)$, without identifying it with the identity.

For two such layers, \eqref{eq:graded-interchanger} has coefficient map
\[
 \chi:P\otimes P\longrightarrow(U\otimes U)\times_k(U\otimes U).
\]
Its source dimension is $9$ and its target dimension $7$. With basis $e,u,v$ of $P$, $e\otimes e$ maps to the shared augmentation coordinate; $e\otimes u,u\otimes e,u\otimes u$ map to the three first-leg coordinates, and $e\otimes v,v\otimes e,v\otimes v$ to the three second-leg coordinates. The two cross terms $u\otimes v,v\otimes u$ map to zero. Thus $\chi$ has rank $7$ and a two-dimensional kernel. It is an actual augmentation-preserving native map over an identity raw carrier map and is not invertible. This detects structure that the effect-only carrier calculation cannot detect.

The constant-singleton effect monad permits empty raw carriers. Neither it nor existence of a coefficient limit certifies physical feasibility or progress. Even a singleton raw carrier can have no normalized coefficient witness: the augmented maps $(k,\id)\rightrightarrows(k^2,\operatorname{pr}_1)$ given by $t\mapsto(t,0)$ and $t\mapsto(t,t)$ have equalizer the zero space, whose augmentation has no preimage of $1$. If a model requires such normalized witnesses, their existence is an additional admission check; augmentation-preserving native maps then transport any supplied normalized witness. Other restrictions on coefficient admission must likewise be preserved by the constructed maps. The displayed nonfaithfulness/interchange examples have explicit normalized vectors. They test native categorical structure, not an external substrate's resource accounting.

\subsection{A sharp obstruction to existence-only lifting}

Coefficient limits are a sufficient constructive hypothesis, not a necessary characterization of every possible native doctrine. Their role is stronger than showing that each comparison happens to have some filler: they construct actual fillers with identity/composition laws. That distinction is forced by the following obstruction.

Define a category $\mathcal K$ with objects $J_0,P_0$, endomorphism groups $C_4$ and $C_2$, respectively, and $\mathcal K(J_0,P_0)=C_2$, with no reverse arrows. Composition within the groups is addition; postcomposition acts on $C_2$ by addition, and precomposition by $a\in C_4$ adds its residue $\bar a\in C_2$. These commuting actions make a category. Fix the leg $\ell=0:J_0\to P_0$. For every primitive comparison $g\in\operatorname{End}(P_0)=C_2$, there are native lifts $a\in\operatorname{End}(J_0)=C_4$ satisfying
\[
 \ell a=g\ell\quad\Longleftrightarrow\quad \bar a=g. \tag{NS}\label{eq:nonsplit}
\]
But an identity- and composition-preserving choice $g\mapsto a_g$ would be a group section of $C_4\twoheadrightarrow C_2$. The nonzero element of $C_2$ has order two; its lifts are odd elements of $C_4$, whose doubles equal $2$, not $0$. No such section exists.

This is an obstruction inside a valid original native presentation. Take the free strict monoidal category on $\mathcal K$, with objects finite words and length-preserving coordinatewise arrows. Singleton-word homs and the equations above remain unchanged. Over the terminal pure category, its constant projection and extension, concatenation accumulation and empty-word unit form a strict contextad. With trivial content, the original comma category is that very native grade category. Every carrier and body map is the identity of $1$, and every individual filler in \eqref{eq:nonsplit} exists, but the required coherent lifting does not.

The obstruction concerns a doctrine exposing a retained joint grade through a non-detecting primitive leg. It does not contradict Theorem~\ref{thm:grade-lifting}, whose joint coefficient objects and comparison maps are constructed by the specified limit property. It proves that replacing that construction by ``a native lift exists for each primitive comparison'' is invalid. In a different doctrine the exact extra requirement is a coherent constructor on the actual primitive comparison category, including its equations; a universal fibration or opfibration of the extension functor is not necessary. Retaining additional lift witnesses in the input can change that comparison category, but must be declared and does not supply a canonical section of the original one.

\paragraph{Consequences and surviving boundary.} The new theorem discharges native grade lifting and its generated coherence for the parameter class above, including the original finite augmented-vector-space context doctrine and independently graded writer content. It supplies actual joint grade arrows before the mixed body square and does not assume a whole-system comparison. The original plant-kernel, information/progress and finite guarded-feedback tests therefore apply after forgetting or interpreting these additional grades, with their previous information, admission, coverage and completion obligations intact. Any semantic observation of coefficient data additionally needs a grade-respecting interpretation; forgetting the coefficients is not evidence that those observations agree. Arbitrary original contextads/contentads, instantaneous feedback, physical admission and infinite-effect completion remain outside the proved class.


\section{Feedback closure of the native grade construction}\label{sec:graded-feedback}

Retained-state closure, finite unfolding, native inverse systems and completed infinite objects are distinct targets. Closure retains a state inside a machine; finite unfolding constructs effects and native grade maps at each horizon; an inverse system also has deletion arrows between horizons. An infinite effect or a single infinite native grade requires further existence and preservation results. This section proves the bridges needed for a seeded, guarded, history-retaining interpretation, without constructing a universal completion.

\subsection{Existing constructions and their precise applicability}

Write $\mathsf{Pure}$ for the pure inclusion, distinct from the content coefficient unit $J$. The sequential part of \emph{Native Intertwiner Lift of State Feedback}, using \emph{Directed State Extension of Concurrent Pure Actions}, applies after verifying the pure actions. For $f\in H(A,B)$, holding a pure set $U$ gives
\[
 L_Uf\in H(U\times A,U\times B),\qquad
 ((u,a),w)\longmapsto((u,b),o)
\]
when $f(a,w)=(b,o)$. Keep the context fibre, effect, coefficient, content grade and native coefficient maps unchanged. Holding through a dependent composite retains the same witness pair, bind effect and coefficient tensors, giving an invertible action compositor. For a pure $k:U\to V$, both sides of
\[
 L_Uf\,;\,\mathsf{Pure}(k\times1_B)\cong \mathsf{Pure}(k\times1_A)\,;\,L_Vf
\]
send $(u,a,w)$ to $((k(u),b),o)$, with the same grade maps after unitors. Action associativity, units, naturality and centrality follow from these coordinate equations and monoidal coherence. No arbitrary coefficient is discarded; neither commutativity of $T$ nor symmetry of $\mathcal V,\mathcal D$ is needed.

Thus the existing native-intertwiner construction supplies machines $(X,f)$, with $f\in H(X\times A,X\times B)$, and actual pairs
\[
 (h:X\to Y,\ \alpha:f\,;\,\mathsf{Pure}(h\times1_B)\Rightarrow \mathsf{Pure}(h\times1_A)\,;\,g).
\]
For an exposed outer state $U$, its retained-closure operator is
\[
 \mathsf B_U(X,f)=(X\times U,f), \tag{RF}\label{eq:retained-feedback}
\]
with pure reassociations; arrows become $(h\times1_U,\alpha)$ with those frames. The existing tightening, unit-state vanishing and nested-closure joining laws therefore retain the full native coefficient maps. This is closure in the state construction, not an instantaneous trace $H(U\times A,U\times B)\to H(A,B)$. It chooses no initial state and does not itself execute a run.

The older independent-parallel state theorem additionally requires an actual parallel functor and pure-layer exactness. Coefficient limits and arbitrary joint couplings do not supply these. In the narrower commutative-$T$, symmetric-$\mathcal V,\mathcal D$ case, independent parallel can be defined by product witnesses, Fubini effects and coefficient tensors. Its independent sequential interchanger is the witness shuffle, with commutative-monad interchange and coefficient symmetries, so it is invertible and supplies pure-layer exactness. The older separated-strength and four-process square then apply to this independent-state operator. They do not invert the shared-state map \eqref{eq:graded-interchanger}.

The existing \emph{Guarded Addressed Unfolding and Native Transport} supplies addressed raw prefixes and a final raw coalgebra for its polynomial-monadic content instance and affine commutative effects. It does not contain the new nonfaithful coefficient maps. For a fixed emitted alphabet $E$, use the free-word monoid $E^*$ and the one-child Cartesian polynomial monad $S(Y)=E^*\times Y$, emitting singleton words. On comparisons preserving that emitted alphabet, after the explicit public projection and forgetting coefficients, this supplies a compatible unary case; it does not prove the full graded result below. Its simultaneous finite-frontier product uses a specified Fubini independence policy, which cannot replace an actual interfering coupling. A graded multi-occurrence interpretation remains a separate crossing/interpretation obligation.

The \emph{Recursive Occurrence Processes and Finite Bridge} already represents unbounded calls, spawn, messages, suspension, interference and infinite raw histories. Its strict model-map theorem concerns operational world/store/outcome data and permitted observers. It does not interpret all these constructors in the present graded $H$. Adding unit coefficients to administrative steps and grades to primitives would decorate histories, but would not itself interpret addressed multi-hole returns or whole-left-phase Bind as this one-hole composition. The operational theorem remains credited at its own scope.

\subsection{Common-state joining and retained closure}

Use a finite connected boundary diagram $f_i\in H(X\times A_i,X\times B_i)$ with boundary endpoint maps $1_X\times a_{ij}$ and $1_X\times b_{ij}$, actual native grade squares and a retained compatible coupling section. Connectedness identifies its endpoint limits with $X\times\lim A_i$ and $X\times\lim B_i$. Alternatively retain $X$ explicitly and restrict the raw/grade construction to its common diagonal before assigning effects. A common returned state is an actual boundary equation.

\begin{theorem}[Common-state native closure bridge]\label{thm:common-state-feedback}
Keep the external interface diagrams $(A_i),(B_i)$ fixed. Suppose the primitive comparison endpoints are $h\times1_{A_i}$ and $h\times1_{B_i}$ for one $h:X\to Y$, and suppose the retained effects satisfy $TR_a(q_a)=q'_{ha}$ at every input. Its admission, and fixed-chooser descent when constructing it, are required. Then the comparison induces a joint native intertwiner with that same $h$. For a further common exposed state $U$,
\[
 \mathsf B_U\mathsf J_X(D,q)
 \cong\mathsf J_{X\times U}\bigl(\mathsf B_U^{\rm common}(D,q)\bigr).
 \tag{CJ}\label{eq:common-closure-join}
\]
Here $\mathsf B_U^{\rm common}$ replaces the single common state by $X\times U$ throughout the diagram. The comparison is native-natural and respects identities, composition, closure units and nested closure joining. The closure square for the forward join/substitution comparison also commutes.
\end{theorem}
\begin{proof}
Removing the pure unit frames turns the component intertwiners into exactly the changing-endpoint cells of Theorem~\ref{thm:grade-lifting}. Its constructed grade maps and body square give the joint intertwiner. Keep $h$ itself; coefficient-limit functoriality proves identity and composition on the resulting pairs.

Both routes in \eqref{eq:common-closure-join} use the same input $(x,u,a)$, matching raw tuples, chosen $q_{x,u,a}$, coefficient diagrams and limits, content limit and returned body. Only state association changes. Identity witness maps and canonical actual limit comparisons give the isomorphism. Testing against primitive projections proves naturality; state reassociation and uniqueness prove the unit and nested-closure equations.

For the directed join/substitution cell, the source forwards the same updated exposed state through both layers. Its coefficient projections are $\pi_i\otimes\rho_i$. Closing that cell or constructing it after common closure gives these same projections, the same retained history map and the same transported admitted effect. Limit uniqueness proves the square, and monoidal coherence proves its iterated laws. No inverse interchanger is used.
\end{proof}

The word \emph{common} is essential. Ordinary composition of independently closed machines retains two state factors, unlike chronological composition of exposed bodies reusing one state. There is no general isomorphism between:
\[
 (\mathsf B_Xf);(\mathsf B_Xg)\qquad\text{and}\qquad\mathsf B_X(f;g).
 \tag{NC}\label{eq:no-independent-fusion}
\]
For bit state, let $f$ flip and $g$ read. From separate initial copies $(s,s)$ the left side returns $s$; shared closure returns $1-s$. The diagonal fails the proposed output square. State cardinalities four and two also preclude an invertible state-map comparison. Deterministic effects and unit coefficients suffice for this counterexample.

The hom-product/synchronization results cannot automatically replace the direct proof. A product of two fair-bit contexts in the native finite-probability category would have one projected joint law. Diagonal and independent coupling cones would both have to factor through it by effect-preserving maps, forcing their different laws equal. Extra raw coordinates do not repair this. Selected native synchronization cones, their creation over state cones and output-action preservation remain additional hypotheses.

\subsection{The missing native truncation law}

Same-horizon native comparisons follow from Theorem~\ref{thm:graded-rules}. For coherent truncation first require each admitted conditional effect to be locally normalized:
\[
 T!(p)=\eta(*)\in T1. \tag{EN}\label{eq:effect-normalization}
\]
Affinity is sufficient but stronger. Normalization is preserved by native pushforward and dependent bind. A nonempty joint inherits it from a normalized marginal because $T!q=T!(Tu_iq)=T!p_i$.

Supply augmentations $\epsilon_C:C\to I$ and $\delta_D:D\to J$ on the admitted context and content coefficients, respectively. On the actual admitted maps and tensor expressions require
\[
 \epsilon_{C'}\alpha=\epsilon_C,\qquad \epsilon_I=\id_I,\qquad
 \epsilon_{C\otimes E}=\lambda_I(\epsilon_C\otimes\epsilon_E),
 \tag{AD}\label{eq:relative-discard}
\]
and the analogous content equations. These are relative conditions, not presumed on every arrow of the ambient categories. They can be defined recursively on free tensor expressions, with the stated associator/unitor compatibility.

\begin{lemma}[Tensor-local deletion]\label{lem:graded-discard}
The maps $d_{C,E}=\rho_C(1_C\otimes\epsilon_E):C\otimes E\to C$ satisfy
\[
 d_{C',E'}(\alpha\otimes\gamma)=\alpha d_{C,E},\qquad
 d_{C,E\otimes F}a_{C,E,F}=d_{C,E}d_{C\otimes E,F}.
 \tag{DT}\label{eq:discard-coherence}
\]
The same holds for content. With \eqref{eq:effect-normalization}, these supply the coefficient and effect parts of coherent prefix deletion.
\end{lemma}
\begin{proof}
Substitute the definitions. The first equation uses augmentation naturality and the unitor square. The second uses multiplicativity: both routes tensor $1_C,\epsilon_E,\epsilon_F$ and then apply the coherent unitors. For effects, discarding the next conditional witness after a fixed prefix $p$ gives $\eta(p)$ by normalization; bind therefore returns exactly the previous prefix effect.
\end{proof}

This is a relative minimality statement for a reusable tensor-local, unital, natural, grouping-compatible deletion constructor. Its one-step-to-zero instance supplies $\epsilon_C$; naturality gives $\epsilon_{C'}\alpha=\epsilon_C$; grouping gives multiplicativity. Naturality along the admitted map $\epsilon_E:E\to I$, together with deletion of a unit, forces $d_{C,E}=\rho_C(1\otimes\epsilon_E)$. The effect-to-zero equation is exactly normalization. An isolated tower might admit ad hoc maps under weaker conditions; that does not establish this reusable constructor.

An independent sufficient construction uses the monoidal slices $\mathcal V/I$ and $\mathcal D/J$, tensoring their augmentation maps. Their units are terminal, and the needed slice limits exist when the requisite ambient limits exist. Reapplying the grade-lifting theorem then constructs the maps; no truncation cell is an input. This adds genuine admissibility data: old coefficients may lack augmentations and old native arrows may not preserve them.

For a nonempty connected augmentation-preserving diagram, an underlying limit inherits $\epsilon_L=\epsilon_i\pi_i$, independent of $i$. Its induced maps preserve augmentation by their projections. So does the forward join/substitution cell: after any target projection its augmentation is the tensor of the source augmentations. For disconnected diagrams use augmentation-preserving cones or slice limits. In vector spaces the ordinary product of two $(\mathbb Q,\id)$ is $\mathbb Q^2$, and no augmentation equals both projections. AugVect already uses the correct slice limits; adding discard must not silently preserve every old joint object.

\paragraph{Sharp coefficient obstructions.} For $T=\mathrm{Id}$ and $\mathcal V=(\operatorname{FinSet},+,\varnothing)$, a singleton step coefficient has no arrow to the zero-prefix unit. Native finite comparisons exist but truncation fails at horizon one. In full $(\operatorname{FinVect}_{\mathbb Q},\otimes,\mathbb Q)$ maps to the unit exist, but unital natural discard on all native arrows is impossible: naturality at $2\id_{\mathbb Q}$ gives $2\id=\id$. Restricting to augmentation-preserving arrows changes the admitted native category.

\subsection{Seeded state closure and full native prefix systems}

After including a strictly delayed controller in retained state, write a body as $F(s,w)=(s^+,b,o)$ with $w\in W_s$, $p_s\in TW_s$ and coefficient $C_{s,w}$. Its fixed writer grade is $Q=(O,(D_o))$. A seed is additional data; comparisons require $s'_0=h(s_0)$. Before controller closure require the native laws at every admitted state/history and command, together with the controller's observation square. The common observation policy is deterministic; a randomized controller is covered only after its own retained innovations, conditional effects and native comparison have been included among the primitives. Retained proof witnesses and coefficient vectors do not become controller observations.

Put $P_0=1$, $P_{n+1}=\sum_{p\in P_n}W_{s(p)}$, with retained bind effects $\theta_n$, and put $G_0=I$, $G_{n+1}(p,w)=G_n(p)\otimes C_{s(p),w}$. Content is $Q^{\boxtimes n}$. For full comma truncation retain output histories
\[
 B_n=S^{n+1}\times B^n.
\]
The body returns the state history, public emissions and ordered writer outcomes. Pure maps $\pi_n:B_{n+1}\to B_n$ and $v_n:O^{n+1}\to O^n$ delete the last entries. These semantic histories do not enlarge the controller's observation.

\begin{theorem}[Seeded graded history-prefix bridge]\label{thm:graded-prefix-feedback}
Under normalization, relative augmentation and admitted joint choices/limits, finite unfolds give actual $F_n\in H(1,B_n)$ and native truncation cells over $\pi_n$, witness deletion and content deletion. They form an inverse system in the total category of mixed operations with fixed input and changing pure output. An object of this category is $(B,F\in H(1,B))$; an arrow to $(B',F')$ is a pure $k:B\to B'$ and an actual native cell $k_*F\to F'$.

All-state native comparisons preserving augmentations, seeds and controller observations give compatible full native maps $R_n:F_n\to F'_n$. On witnesses and effects,
\[
 \tau'_nR_{n+1}=R_n\tau_n,\quad TR_n\theta_n=\theta'_n,\quad
 T\tau_n\theta_{n+1}=\theta_n.
 \tag{PT}\label{eq:graded-prefix-laws}
\]
The square also commutes on actual context/content coefficients and history bodies. Identity/composition, single-wire and nested-wire delayed closure, common-state joining and the admitted augmentation-preserving directed interchanger are preserved.
\end{theorem}
\begin{proof}
Sequentially fold the actual body, forwarding state and recording old emissions by pure history frames. Induction identifies the resulting witnesses, effects, coefficients and contents with $P_n,\theta_n,G_n,Q^{\boxtimes n}$. This is the finite interpretation bridge from the seeded machine, not a terminal-result quotient.

Truncation uses $d_{G_n(p),C_{s(p),w}}$ on context coefficients and deletes the last writer factor using $\delta_{D_{o_n}}$. Lemma~\ref{lem:graded-discard} and normalization prove the native grade/effect equations. Pure history deletion and writer-outcome deletion prove the full body square. Block deletion follows from \eqref{eq:discard-coherence}.

At matched histories the controller selects the same command. Set $R_{n+1}(p,w)=(R_np,r_{s(p)}w)$, tensoring the prefix coefficient map with the actual local $\alpha$; content maps tensor the actual $\beta$ maps. The conditional body square identifies the target input. Bind gives effect transport, augmentation preservation gives the truncation square, and coordinatewise state maps commute with history deletion. Identity and composition follow inductively.

A delayed wire enlarges state to $(s,d)$, uses the old $d$ and stores the returned $d^+$ for the next step. Witnesses, effects and coefficient maps are those at that actual input. Thus unfolding the retained closed machine equals this graded construction on the enlarged state. Two wires use the same old pair and new pair under both closing orders; reassociation and paired seeds give the laws. The common-state theorem identifies the one-step data on both sides of \eqref{eq:common-closure-join}; induction gives all finite unfolds. Its limit maps and directed interchanger preserve augmentation, so these comparisons also commute with truncation.
\end{proof}

Terminal-state output alone need not admit full comma truncation. A transition sending both $0,1$ to $0$ has no pure output map recovering both predecessors; with a fixed seed, use a random first step and then this reset. Context truncation can exist while the output-deletion field fails. History outputs supply that field. The older Set/effect-prefix theorem and same-horizon body squares therefore cannot alone be read as a full graded $H$-tower theorem.

The selected unbounded semantics is the native inverse system together with compatible infinite raw histories. The raw comparison acts coordinatewise; reverse path lifting needs separate local coverage. No inverse limit inside $H$, instantaneous solution selector or yanking law is asserted. Fixed-open-input law preservation still does not replace the all-history conditional hypotheses, as the existing null-diagonal counterexample shows.

\subsection{Separate completion boundaries}

Compatible fair-bit laws need not yield an infinite $\mathcal D_{\rm fin}$ effect, as already proved. Grades have an independent obstruction even with $T=\mathrm{Id}$ and singleton raw histories. In AugVect repeat $U=(\mathbb Q^2,\operatorname{pr}_1)$, obtaining $U^{\otimes n}$ with deletion $1\otimes\epsilon$. Put $e=(1,0)$ and $u=(0,1)$. For each position $m$, the perturbation zero before $m$ and thereafter equal to a tensor with $u$ at $m$ and $e$ elsewhere is compatible. Adding it to $e^{\otimes n}$ gives a cone from the augmented unit.

If a finite-dimensional augmented limit existed, the vectors representing these cones minus the base cone would be infinitely linearly independent: projecting to a stage containing any finite collection of positions separates their tensor-basis coordinates. This is impossible. Native inverse systems therefore need not have one infinite coefficient object, even when the raw limit and effect are trivial. Content has the same possible obstruction.

The conditional \emph{Continuous Completion of Display Interaction} theorem applies only after verifying limits, preservation and typed transport in the actual native presentation. It is not automatic for this coefficient model; finite distributions and finite-dimensional coefficients fail relevant conditions. No universal completion or completed infinite measure is used for the raw-path and finite-cylinder judgments below.

\subsection{A combined stochastic, graded, interfering recurrent realization}\label{sec:combined-feedback-model}

Take $T=\mathcal D_{\rm fin}$ and $\mathcal V=\mathcal D=\operatorname{AugVect}_{\mathbb Q}$. Put $U=(\mathbb Q^2,\operatorname{pr}_1)$ and $P=U\times_{\mathbb Q}U\cong\mathbb Q^3$. A configuration has one shared $x\in\mathbb N$, two worker phases and read registers, retained sample bits and a control phase. Histories retain all primitive witnesses, scheduler choices and intermediate states. Scheduler observations exclude the source-only auxiliary bits below. The delayed controller observes the completed counter.

An epoch has six microsteps: joint sample, four scheduled read/write actions, and restart. Sampling draws $(b_1,b_2)\in\{0,1\}^2$ under the retained diagonal, antidiagonal or independent fair-bit coupling $q$. Its raw carrier always has all four pairs, including zero-weight cases. Component coefficients are $U$ and boundary maps are augmentation, so the joint coefficient is $P$. Worker $i$ then executes
\[
 R_i:r_i:=x,\qquad W_i:x:=r_i+1+b_i.
\]
All six read-before-write schedules are admitted. Each worker step selects an enabled action; this particular model has no extra idle step or external mutation. Other microsteps have context coefficient $I$. Restart clears the completed worker registers, selects the next coupling from $x\bmod3$ (diagonal, antidiagonal, independent respectively), and returns to sampling. The sample uses that previously selected command. This is guarded recurrent state passing, not instantaneous solution of an equation.

Use one fixed tagged writer grade for the dispatcher. Sample-bit-pair tags carry coefficient $P$; read/write/restart tags carry $I$. The body emits its actual tag and returns the actual next configuration. Values and state changes remain in full boundary histories. After $m$ samples, the context and realized content coefficients are $P^{\otimes m}$. Truncation discards an actual sample factor through augmentation, or a unit factor at another microstep.

In the source interpretation each component also samples an auxiliary fair bit $a_i$, with
\[
 q^{\rm src}(b_1,a_1,b_2,a_2)=q(b_1,b_2)/4.
\]
The target erases these bits and their private records. Since scheduler and controller observations exclude them, the next coupling chooser descends through the state projection. Give the first component coefficient comparison identity and the second $s=\operatorname{diag}(1,0)$. The constructed joint context and sample-content map is
\[
 c:P\to P,\qquad c(a,u,v)=(a,u,0). \tag{CM}\label{eq:combined-coefficient-map}
\]
It preserves augmentation and is noninvertible. Choosing identity instead gives a distinct native comparison over the same raw map; the normalized vector $(1,0,1)$ separates them. Thus nonfaithful grades and nontrivial probability occur in the same recurrent system.

Every rule has an explicit core interpretation. The sampling diagram shares the world/control boundary and keeps the two workers' registers as distinct component interfaces; their limit supplies the one configuration. No component sample is declared to update a duplicated copy of the entire configuration. Joint sampling uses Theorem~\ref{thm:grade-lifting} on the local sample grades and their common boundary. Reads and writes are deterministic primitive bodies on one counter, framed by the other worker's local state. Cases select commands and phases. Conditional bind retains the registers and sampled bits. Restart updates state for the next guarded step. The state projection, local effect maps and \eqref{eq:combined-coefficient-map} satisfy the conditional native equations at every admitted raw input and command. Theorems~\ref{thm:graded-rules}, \ref{thm:common-state-feedback} and \ref{thm:graded-prefix-feedback} therefore supply one comparison for every compatible causal scheduler and every horizon. No whole-system kernel equation is an assumption.

Dependence has a consequential effect. Serial execution increments by $2+b_1+b_2$: diagonal gives $2,4$ with masses $1/2,1/2$, antidiagonal gives $3$ surely, and independence gives $2,3,4$ with masses $1/4,1/2,1/4$. Under $R_1R_2W_1W_2$, the last write instead gives only $1+b_2$. Bits $00$ therefore produce $x+1$, whereas either serial order gives $x+2$. The changed counter controls subsequent coupling choices. Correlation, interference and feedback thus act on the same histories. A stale write can decrease the counter relative to an earlier microstep; those intermediate effects are retained.

There is a scoped universal progress result. Each enabled worker action decreases the four-action remainder by one; completed workers cease to be eligible. An invariant during the epoch is that the counter and every initialized read register are at least the epoch's initial counter. Reads preserve it, and writes exceed that initial value because $1+b_i\ge1$. Hence the completed counter increases by at least one. Sampling and restart each take one step, so every six steps complete an epoch and $x$ after epoch $m$ is at least $x_0+m$. Every maximal raw run is infinite in this total no-idle model, including all zero-weight sample cases. The comparison preserves phases and $x$, so it preserves this progress guarantee. This is step productivity under the explicit rules, not a fairness, elapsed-time, resource-capacity or arbitrary-environment theorem.

The exact diagnostic checks three initial couplings and every pair of the six schedules over two epochs, including adaptive next-coupling selection. It compares step closure with an independently written full-choice fold, verifies prefix law pushforward and effect truncation, and retains zero-weight raw paths. Matrix checks distinguish the native maps and verify the augmentation truncation square on both grade sides. The tests corroborate the uniform written induction; they do not define a bounded semantics. The selected unbounded semantics consists of raw histories, consistent finite cylinders and compatible native grade prefixes, with no infinite $\mathcal D_{\rm fin}$ effect or finite-dimensional infinite tensor.

\section{A second realization: suspended asynchronous occurrence processes}\label{sec:occurrence-realization}

This realization uses the existing occurrence machine rather than a second parameter choice of the counter model. Its distinguishing structure is asynchronous FIFO messages with fresh identities, blocked receive, resumption after interference, and recursive one-hole continuations with actual internal jobs. We interpret the smallest fragment needed for these features, not the whole occurrence language. Context and content coefficients are terminal; the native comparisons still retain noninjective witness maps and actual conditional effects.

\subsection{Selected syntax, configuration and observation}

Use $\mathsf{Ret},\mathsf{Op},\mathsf{Bind},\mathsf{Call},\mathsf{Send},\mathsf{Receive}$, with one $\mathsf{Spawn}$ during initialization. Each completed operation/result tree has exactly one addressed hole; its surrounding event nodes remain retained. The restriction is inductive: return/send/receive/spawn have one hole, admitted operation tables have one, and grafting one one-hole continuation into a one-hole left result has one. Calls use a finite table of such definitions. Divergence is permitted and is not assigned a fictitious completed hole. No $\mathsf{Batch}$, multi-hole $\mathsf{Await}$ or authority transfer is needed. Capabilities and the frontier token ledger are empty in this family. This removes resource obligations, not shared mutable state or message identity.

All ordinary data functions and continuation/procedure families are deterministic, total, terminating mathematical functions of their declared data. They do not inspect the machine, hidden audit state or future random outcomes. The Python callbacks realize this stipulated interface; the runtime is not a sandbox enforcing it.

The interpreted state $\Sigma$ is the actual finite occurrence configuration: shared world, every job's identity/rank/term/value/mode, parent and continuation/join records, completed results and claim flags, result trees and their raw provenance, queues, messages, both fresh allocators, primitive events and administrative trace. Histories additionally retain each selected command, primitive row or message identity, and the intermediate configuration. Two public roles do not mean two runtime jobs: recursive Bind leaves finite but unbounded families of internal children and waiting ancestors. None is erased.

The chosen queue interpretation is reliable FIFO: Send appends, and Receive selects index zero. The runtime exposes other message indices; FIFO is our admitted policy, not a theorem about every runtime call. A Receive remains in evaluation mode while its queue is empty and is absent from the enabled menu. Send and Receive are not both recorded in the primitive event word, so that word is not the history observer: use the complete step/configuration history and raw result records. A scheduler selects an enabled command whenever the global menu is nonempty. Stopping an observer at a fuel limit, returning no scheduling choice prematurely, or using a root-completion stop flag is not global quiescence.

\subsection{A concrete suspended stale-update protocol}

The public world contains one shared bit $g$ and a completed-epoch counter $c$, initially zero in the diagnostic. Define a client and a server, with dedicated request and acknowledgement channels:
\[
\begin{aligned}
C(n):\quad&
r\leftarrow\mathsf{read}(g);\
\mathsf{send}(\mathrm{req},(n,r));\
a\leftarrow\mathsf{receive}(\mathrm{ack});\
\mathsf{commit}(n,r,a);\
C(n+1),\\
S:\quad&
(n,r)\leftarrow\mathsf{receive}(\mathrm{req});\
b\leftarrow\mathsf{flip}(g);\
\mathsf{send}(\mathrm{ack},(n,b));\
S .
\end{aligned}
\]
Semicolons here abbreviate the actual one-hole Bind syntax, not atomic native composition. The commit writes $g:=1-r$ using the saved ordinary-data capture; the acknowledgement's epoch is checked in the declared protocol domain. In this closed two-role protocol the request/acknowledgement invariant supplies that domain. Broader clients injecting malformed payloads require their own explicit refusal rules, not an unmodelled Python exception. The declared read input is $n=c$ and it returns $(n,g)$. Flip takes $(n,r)$ with $n=c$ and returns $(n,1-g)$ while storing $1-g$. Commit takes $(n,r,(n,b))$ with $n=c$, stores $g:=1-r$ and $c:=n+1$, and returns $n+1$. The schematic notation above suppresses the unchanged epoch tag on read/flip results. The fresh message identity returned by Send is not confused with its payload. The client keeps $r$ in a continuation capture, not in the dynamic Send result.

The primitive domains are invariant in the closed protocol: client $C(n)$ starts with $c=n$; sending and server flip do not change $c$; the matching acknowledgement precedes the sole commit that changes $c$ to $n+1$; only then can the client send the next request. Unique claims and FIFO identity prevent duplicate execution of that continuation. This argument is independent of the audit bit. The diagnostic callback assertions check these domain premises; a failed host assertion is not a semantic refusal. Admitted additional contexts must preserve the primitive and payload domains, or supply explicit preserved refusal branches before they can be included.

Initialize by spawning the server with an empty bundle and continuing with $C(0)$, also with an empty bundle. Retain every bootstrap and continuation allocation step. Recursion unfolds in an actual Call transition; each protocol cycle also contains communication and primitive transitions. Unbounded calls do not require a finite rank or a fuel limit.

Starting with $g=0$, the client can read $0$ and suspend awaiting an acknowledgement. The server receives its request, flips the same bit to $1$ and enqueues the acknowledgement. The resumed client commits from its stale $r=0$, leaving $g=1$. In contrast, two serialized read/modify/write toggles return the bit to $0$ in either order. Moreover neither complete request/acknowledgement method can execute atomically from the initial empty queues: each requires the other's communication. This is a decisive failure of whole-method atomic/serial abstraction. A rendezvous model could still produce the same final bit. Its separate boundary is the declared history observer: after asynchronous Send completes but before Receive, the pending message and resumed sender are actual retained state. A synchronous model that erases that interval does not preserve this observer without a separately justified quotient.

The example is a family over initial bit and epoch, finite procedure/continuation parameters of this form and admitted public schedulers. A finite recurrent schedule constructs arbitrarily long executions by repeating the handshake with increasing epoch; each configuration remains finite while its allocator/control history grows. We do not require a fairness or physical-time theorem to interpret these runs. Arbitrary fragment programs can terminate, block or unfold forever, and those cases are retained below.

\subsection{Local interpretation into the original presentation}

Take $T=\mathcal D_{\rm fin}$ and $\mathcal V=\mathcal D=\mathbf1$ in Proposition~\ref{prop:graded-presentation}. An enabled rule at command $\ell$ has an input domain $\Gamma_\ell=\{\sigma:\ell\in\operatorname{Enabled}(\sigma)\}$, a raw witness family $W_{\sigma,\ell}$, its conditional probability $p_{\sigma,\ell}$ and a body returning the updated configuration and a retained rule/event tag. This gives an original mixed object in $H(\Gamma_\ell,\Sigma)$ with writer outcomes. Construct these objects from the following local data, before forming a dispatcher.

\paragraph{Operation rules.} For each declared read/flip/commit input, the primitive table supplies its actual row $z$, probability, public result/event and next local world. Its job-result view uses that same row to construct the one-hole result tree and raw record. A shared transaction boundary retains the operation, alternative, row identity and public result. The world and job views are joined over this boundary with the \emph{same} chosen primitive effect. The matching carrier is the graph of their common row; its joint effect is the primitive law pushed onto that graph. Independent resampling of world and result is not allowed. The frame preserves all other jobs, queues and allocator fields. Only the primitive operation table is supplied, not a compound program kernel.

\paragraph{Queue rules.} For Send, use a job view, a channel/message view and an allocator view. Their shared transaction boundary is
\[
 e=(\mathrm{send},\mathrm{pid},\mathrm{mid},\mathrm{channel},\mathrm{payload}),
\]
with $\mathrm{mid}$ the current fresh-message counter. One view finishes the job with that ID, one appends that exact message and payload, and one advances the allocator. Their input/output overlap equations and Dirac witness effects give a local star-shaped join. For Receive, the guard is a currently nonempty queue. Its boundary identifies the actual head message and payload; the queue view removes that one ID and the job view returns that payload. The message is retained between Send and Receive, which are separate transitions. The join does not synchronize an entire sender with an entire receiver or remove the suspension interval. Empty token transfers are explicit in this interpretation; no physical ownership claim is inferred.

\paragraph{Control rules.} Return constructs its actual one-hole result. Bind allocates the left child and changes the parent to its waiting-left mode. When that child completes, the separate parent step claims it, extracts its unique actual leaf, allocates one continuation child with $k(v)$ at that value and post-interface, and records the parent join. A later step claims the completed continuation and grafts its actual tree at the retained address. Call evaluates the public procedure argument, unfolds its declared body and emits the administrative unfold record. The initializing Spawn allocates one public server job and returns its handle. These are pure conditional updates joined with the relevant job/allocator views and framed by untouched state. They have singleton witnesses, Dirac effects and unit coefficients; they are not suppressed transitions.

The local joins have independently specified state/effect boundaries, so Theorem~\ref{thm:grade-lifting} constructs their native maps and full bodies. Their raw records identify one transaction, rather than duplicating it into several independently counted events. Typed enabled-command cases and framing assemble the dispatcher. Its next-state law is consequently a derived conditional law. FIFO identity, freshness, claim flags and dependent continuation inputs are concrete local equations; none is an admission condition saying that the final whole program already compares correctly.

\paragraph{Dependent continuation and phase compatibility.} The map
\[
 (\theta,v,\sigma_{\mathrm{current}})
 \longmapsto
 \bigl(\text{fresh child at the unique address of }\theta,\ k(v),\
        v,\sigma_{\mathrm{current}}\bigr)
\]
uses the completed left tree and the current shared state, not the state at Bind entry. Native comparison preserves the public leaf, address, allocator and continuation family; therefore this allocation square commutes. Claims and grafts commute by their explicit tree and flag maps. Chronological composition of these rule objects is exactly the core's dependent substitution, retaining every intervening configuration and witness. Syntactic Bind elaborates to this control protocol; it is not identified with one atomic $H$-composite.

That distinction is necessary already with one hole. In $\mathsf{Bind}(\mathsf{Ret}(*),k)$ where $k$ reads a shared bit, another job may write the bit after the left return and before continuation execution. Atomically executing the bare return/read composite at Bind entry forbids that result and also omits allocations and scheduling records. The smallest justified repair is the retained control-frame elaboration just given. No new multi-hole content mechanism is needed for this fragment. Conversely, the proof does not identify different source Bind bracketings at the full chronological/identity observer merely because core chronological microstep composition is associative.

\subsection{Exact operational interpretation, including stopping}

Let $\operatorname{Run}_{\rm occ}(t)$ be the prefix tree of admitted occurrence-machine derivations of a term in this fragment, with scheduler commands and raw rows retained. Let $\operatorname{Elab}(t)$ be the initialized native state machine assembled from the above rule objects, with state the actual $\Sigma$ and the same enabled command domain.

\begin{theorem}[One-hole occurrence interpretation]\label{thm:occurrence-interpretation}
For each primitive interpretation, finite occurrence derivations are in bijection with the well-typed finite native microstep derivations of $\operatorname{Elab}(t)$. This preserves complete configurations, addressed one-hole trees, job/message identities, primitive and administrative chronology, conditional weights, and truncation. Compatible infinite real histories correspond coordinatewise. Local suspension, global quiescence and actual Call divergence are distinguished.
\end{theorem}
\begin{proof}
A labelled occurrence step, including its selected raw row, determines its core witness uniquely: its primitive row, head message, or singleton administrative witness. The local transaction joins have exactly the matching records described by the runtime rule. Their body updates the same fields, and the frame leaves precisely the other fields unchanged. Bind has the three distinct allocation/claim/graft phases specified above, so this case uses the whole-left completion and the actual continuation value. Fresh counters and flags determine the same identities and readiness. These cases give inverse one-step translations; none assumes equality of a preassembled global kernel.

Induct on the number of real steps, using the actual next configuration to type the next witness. Products of the same conditional probabilities agree, without treating scheduler alternatives as independent random samples. The inverse translations delete the last step in the same way. Thus a compatible infinite derivation maps by its finite prefixes and has the inverse coordinate map. An empty queue prevents the same Receive transition on both sides, while a Call-to-Call loop still supplies infinitely many actual unfold transitions.
\end{proof}

The normalized all-horizon prefix theorem cannot be applied to an empty enabled domain without an explicit interpretation. Use the occurrence source's stopped-history convention when necessary. Classify a configuration as live if its declared global command menu is nonempty, otherwise as $\operatorname{Stop}(\kappa,\sigma)$, retaining terminal/deadlocked status and the complete configuration. Local blocked receive does not make the whole machine stopped while another job is ready. An observation cutoff or root-only stop is not this classification.

The semantic wrapper gives a stopped state one Dirac, unit-grade transition labelled $\operatorname{Pad}(\kappa)$. This is not a runtime transition. A horizon-$n$ operational observation is a history of exactly $n$ real steps, or a maximal stopped history of length strictly less than $n$; fill the latter's unique marked padding suffix. This is a bijection with the wrapper's native $n$-prefixes. Truncation deletes a real last step or leaves an already-stopped record fixed. At infinity, decode an eventually padded history to its finite maximal real history; a never-padded history is an infinite real execution. Therefore padding is not mistaken for physical divergence.

This wrapper uses only typed pure cases and unit primitives already available in $H$. Its normalization and coefficient deletion hold because finite distributions are normalized and both coefficient categories are terminal. Theorem~\ref{thm:graded-prefix-feedback} constructs its full history-output prefix system. For a fixed compatible causal scheduler/strategy $\pi$, the resulting effect-bearing interpretation is
\[
 \operatorname{OccurrenceModels}^{\pi}_{\rm unary,FIFO}
 \xrightarrow{\operatorname{Elab}}
 \operatorname{InitializedSystems}_H
 \xrightarrow{\operatorname{Unfold}}
 \operatorname{NativePrefixSystems},
\]
with the preceding physical-history decoding. A history-dependent strategy retains its permitted history as deterministic memory; this adds neither physical transitions nor new observation capabilities. A randomized scheduler supplies its innovations and their native comparison among the primitives. These choices supply the conditional finite laws. Demonic commands are retained, not silently mixed into one probability distribution. No infinite $\mathcal D_{\rm fin}$ element is required.

\subsection{Native replacement and actual theorem reuse}

The source implementation stores a private audit bit alongside the public shared bit. Each read, flip and commit can have two fair raw auxiliary outcomes (and an explicitly retained zero-weight row in the diagnostic), all with the same declared public result and public state update. The target has one corresponding raw row with weight one and omits the audit field. Let $h$ erase only the audit state and let $\rho$ map each source row to that target row. Then
\[
 \sum_{\rho z=z'}p(z)=p'(z'),\qquad
 h(\operatorname{post}(\sigma,z))
 =\operatorname{post}'(h\sigma,\rho z)
\]
hold for each actual primitive input. Primitive public event/result trees agree; send/receive and control operations introduce no auxiliary dependence. Costs can be fixed identically if observed. No resource or time interpretation is needed.

On configurations, preserve jobs, terms, values, ranks, allocators, messages, queue order, addresses and claim flags. Map the world by $h$ and recursively map primitive raw records within pending/completed results. Event-row annotations map to the corresponding target table annotation. Source row weights are not equated pointwise with target weights: their required relation is fibre summation, and the interpretation retains each model's actual conditional effects. Public observers must factor through this map. In particular the scheduler cannot inspect the erased row/audit data, and callbacks cannot capture it.

Each local queue/control view square commutes: public values and fresh IDs agree, FIFO heads are identical, and readiness/claim guards have the same truth value. Operation view squares are the displayed primitive equations. The source coupling on shared primitive-row views pushes to the target coupling; administrative joins are Dirac. Consequently Theorems~\ref{thm:grade-lifting} and \ref{thm:graded-rules}, already used for the first realization, construct the dispatcher comparison. We do not posit its global kernel square.

\begin{corollary}[Reused recursive contextual model comparison]\label{cor:occurrence-reuse}
For every well-formed term/procedure table in this fragment, and every admitted fragment context with the stated local primitives, preserved primitive/payload domains, FIFO policy and observation restriction, the replacement above induces compatible full native prefix comparisons. On decoded histories it preserves actual administrative/communication chronology, shared-state projection, finite status and infinite real divergence; finite conditional laws push forward. The construction respects the local joining and dependent substitution used in elaboration and the guarded recurrent unfolding. Universal raw properties of the target transfer to the source along this total history map.
\end{corollary}
\begin{proof}
The local certificates just checked are precisely the hypotheses of the existing generated native-comparison theorem. Normalization and terminal coefficient augmentations discharge the additional hypotheses of Theorem~\ref{thm:graded-prefix-feedback}; seeds and controller/scheduler observations are compatible. Apply those theorems, rather than prove a new system-level comparison induction. The interpretation bijection of Theorem~\ref{thm:occurrence-interpretation} transports their conclusion to occurrence histories.

Readiness and status reflection make stopped padding map to padding of the same kind, so physical-history decoding commutes with the comparison. An infinite real execution cannot become a finite stopped one. Call unfolding is a real one-step rule on both sides; no termination induction is used. Context substitution adds only the same interpreted local rules and their checked domain and observation interfaces, so the same generated theorem applies to the enlarged fragment context. This is a contextual implementation/model-replacement result, not arbitrary equivalence of distinct source terms or of all occurrence-language contexts.
\end{proof}

Admitted interference is explicit. Server writes while the client is suspended, and every other job's primitive acts on the one current world. If a context adds a world-bit environment transition, it must satisfy the same local square and be included in the global command menu; no queue/ledger mutation or new observer is admitted by implication. Readiness is computed after the actual interference. The retained request, acknowledgement and cached read remain in the configuration until consumed; resumption does not restore a snapshot or erase a pending effect.

This is the same nontrivial theorem consequence reused in two different interaction structures. The first realization combines selected stochastic couplings with nonfaithful grades and a shared-counter scheduler. The second uses FIFO storage, blocked resumption and growing occurrence-control records, with terminal coefficients and a directed stochastic row erasure. New work consists of the queue/control interpretation and its local certificates. The generated joint comparison, conditional-law/prefix transport, native truncation and feedback arguments are reused unchanged. The change in primitive implementations is consequential to raw witnesses and conditional probabilities; only the justified observer cannot distinguish it.

\paragraph{Diagnostics and limits.} The accompanying check executes the actual occurrence runtime and an independently written defunctionalized rule interpreter. It compares every retained intermediate configuration and administrative phase, rather than only result values or the primitive event word. It tests the request/acknowledgement loops under several public schedules, source auxiliary erasure and zero-weight rows, plus a two-message FIFO fixture and stopped/divergent boundary cases. These finite checks corroborate the written rule induction; they do not bound the interpretation's execution length or establish a physical channel, host-language security or resource/progress theorem. Full multi-hole returns, Batch, general Await and nonempty ownership transfer remain outside this fragment and require their own interpretation.


\subsection{Compatibility verdicts and contract v0.1.3}\label{sec:core-completion-assessment}

This assessment uses the proposed \emph{Feedback and Concurrency Research Contract v0.1.3}, especially its three completion ledgers. Mathematical support, review and institutional acceptance remain distinct. Verdicts concern the inspected source statements and the bridges proved here, not an exhaustive audit of every repository artifact.

\begin{description}
\item[Established construction.] Coherent pure actions make the existing sequential native state construction and retained closure available. Common-state joining and closure commute by \eqref{eq:common-closure-join}; the closure/interchanger square is directed. Full native history-prefix feedback follows under local normalization, relative augmentations, compatible seeds/observations and actual output deletion. The combined realization verifies these hypotheses simultaneously with probability, nonfaithful grades and shared-state interference.
\item[Additional hypotheses.] Full independent-parallel state laws require their parallel and pure-layer-exactness inputs; they hold in the specified commutative/symmetric independent-state construction, not for arbitrary shared couplings. Native synchronization needs its selected-cone and output-preservation hypotheses. Branching occurrences need a graded occurrence/crossing interpretation. Infinite effects and infinite coefficient/content objects need separate completion and transport results.
\item[Failed unrestricted laws.] Separate closure does not merge state copies. Finite preservation does not imply native truncation without normalization/discard and an output deletion map. Fixed-open-input preservation does not imply adaptive closure. Compatible finite effects and grades do not imply infinite completion. The counterexamples identify the failed fields.
\end{description}

For the categorical clauses, RC2 has a positive independent context/content family and native grade-lifting construction. The construction portions of RC3--RC4 now have retained-state feedback, a seeded guarded native-prefix operator and its common-state interaction laws. The native-comparison portion of RC5 has functorial preservation through these constructors, including raw-path and conditional finite-law consequences. The structural portion of RC7 has the original comma presentation and exact augmented-vector-space recovery. These written construction/coherence results hold at the declared parameterized scope; unfinished language work does not negate them.

The second realization in Section~\ref{sec:occurrence-realization} now supplies an explicit interpretation of unary asynchronous occurrence processes, including suspended receive, local queue/control joining, dependent continuation phases and recurrent comparison. Its exact operational/core history bridge and Corollary~\ref{cor:occurrence-reuse} reuse the same generated native-comparison and history-prefix theorem as the stochastic/graded realization. This closes the previously identified second-interpretation gap for the selected scope; it does not assert an interpretation of the entire occurrence language.

Full multi-hole returns, Batch, general Await and nonempty ownership transfer remain realization-programme work. The existing full runtime's operational contextual theorem and old phase-evaluator bridge retain their own credit and scope. Grade-sensitive observations still need their declared interpretation; instantaneous interconnection and unrestricted return-based recursion are not supplied by the selected guarded semantics. These limitations delimit the established core rather than requiring that every model or language constructor admit its operations.

RC6 receives one combined raw step-productivity and comparison example with retained finite probabilistic dependence. The contract still requires its substantive resource/reservation interpretation and broader selected probability/progress interpretation bridges. Linear coefficients are not physical stock, and a six-step model result is not a general scheduler/deadline theorem. The full realization and application portions of RC1 and RC3--RC8 remain undischarged here.

The core completion clause requires RC2--RC5's positive parameterized construction/coherence/comparison, its RC7 relationship to the original construction, and two materially different realizations verifying the hypotheses and reusing a nontrivial consequence. Those fields now have witnesses at the declared scope: the coefficient-limit original presentation and native joint/rule laws; retained-state and augmented history-prefix feedback; generated contextual model comparison; exact structural recovery; and the stochastic/graded shared-state family together with the asynchronous occurrence-fragment family. The latter's queues, suspension, fresh identities and growing control records materially change its interaction structure. Terminal coefficients satisfy the theorem's hypotheses. This is theorem reuse through a new local interpretation, not merely a second parameter setting of the counter example.

Accordingly, the \textbf{two-realization requirement and the categorical-core completion criterion are met at this declared parameterized scope, as a mathematical-evidence assessment}. This is not institutional acceptance or a claim that every contextad, feedback notion or occurrence constructor has been covered. Requiring the excluded full language before crediting this core would conflate the contract's levels. The realization and application/reuse ledgers remain separately incomplete: the broader required language, resource/progress interpretations and application-comparison evidence have not all been discharged here. No universal infinite-grade completion or downstream Byzantine mission is added as a gate. The contract is unchanged.

\section{Evidence and reusable obligation}

The written proofs establish the effect-fibre fragment and the native coefficient-limit grade-lifting theorem of Section~\ref{sec:grade-lifting}. Exact finite diagnostics in \texttt{checks/check\_local\_joint\_comparison.py} test the coupling/chooser and coverage obstructions, directed maps, shared-state interference and conditional prefix transport. They are not mechanized proofs or a validation of an external substrate. The new exact matrix diagnostics in \texttt{checks/check\_native\_grade\_lifting.py} test native coefficient pullbacks, actual map composition, tensor accumulation, the noninvertible interchanger and the coherence obstruction. The combined feedback diagnostic in \texttt{checks/check\_graded\_feedback.py} compares an independently written fold with step closure, retains zero-weight paths and checks effect/grade truncation and shared-state obstructions. The occurrence diagnostic in \texttt{checks/check\_occurrence\_realization.py} compares actual runtime steps with an independent restricted interpreter and tests native replacement across administrative, queue and primitive boundaries. Source checks do not establish successful typesetting; no proof-assistant certification is claimed.

The reusable obligation for a new substrate is consequently precise: expose actual local effects and boundaries; provide native primitive comparisons; prove local admission transport and, when needed, chooser descent; and interpret the allowed interference constructors with full body/grade squares. The generated kernel square and finite guarded comparison then follow. For the coefficient-limit parameter class, native lifting and constructor coherence are proved here, including the original augmented-vector-space specialization. Beyond that class, native lifting needs a doctrine-specific construction; mere filler existence is refuted by the coherence obstruction above. This retained residual is narrower than assuming the enlarged system is already related, and does not silently replace general Byzantine control by successful examples.

\paragraph{Source boundaries.} The proofs use the local definitions and arguments in \texttt{Tree\_Overlap\_Admission\_and\_Retained\_Reassociation.tex}, \texttt{Native\_Coupling\_Cones\_and\_Conditional\_Assembly.tex}, \texttt{Boundary\_Admitted\_Joint\_Occurrence\_Unfolding.tex}, \texttt{Retained\_Fibre\_Effects\_and\_Guarded\_Prefixes.tex}, and \texttt{Information\_and\_Progress\_Native\_Comparisons.tex}. Their relevant native/conditional statements were inspected for this derivation. This does not independently certify every theorem in those manuscripts. The rule-format approach has established precedent in Turi and Plotkin, \emph{Towards a Mathematical Operational Semantics}, LICS 1997 (primary abstract verified [\(\checkmark\)] at \url{https://www.dcs.ed.ac.uk/home/dt/towards.html}); that work describes natural operational rules inducing distributive laws. No theorem from it is used here to infer the native-grade, shared-world or progress obligations.
\end{document}
